\ifdefined\learninglargeprint
  \documentclass[14pt,oneside]{extbook}
\else
  \documentclass[11pt,oneside]{book}
\fi
\usepackage[T1]{fontenc}
\usepackage[utf8]{inputenc}
\usepackage{lmodern,microtype}
\usepackage[a4paper,margin=27mm,headheight=15pt]{geometry}
\usepackage{amsmath,amssymb,amsthm,mathtools}
\usepackage[most]{tcolorbox}
\usepackage{booktabs,tabularx,enumitem,tikz,needspace,titlesec}
\usetikzlibrary{arrows.meta,calc,positioning}
\usepackage{natbib,xurl,fancyhdr}
\usepackage[hidelinks,unicode]{hyperref}
\definecolor{learnink}{HTML}{18212B}
\definecolor{learnnavy}{HTML}{17365D}
\definecolor{learnpaper}{HTML}{F2F5F8}
\tcbset{colback=learnpaper,colframe=learnnavy,boxrule=0pt,leftrule=1pt,
  sharp corners,breakable,left=3mm,right=2mm,top=2mm,bottom=2mm,
  fonttitle=\sffamily\bfseries,before skip=10pt,after skip=10pt}
\newtcolorbox{learningcard}[2][]{title={#2},#1}
\newcommand{\lessoncue}[2]{\par\needspace{4\baselineskip}\smallskip\noindent\textsf{\textbf{#1.}} #2\par}
\newcommand{\lessonreference}[3]{\begin{tcolorbox}[colback=white,breakable=false,
  colframe=learnink,leftrule=0pt,toprule=0.7pt,bottomrule=0.7pt,
  title={Quick reference}]
  \textbf{Use.} #1\par\textbf{Watch.} #2\par\textbf{Link.} #3
  \end{tcolorbox}}
\newcommand{\claimstatus}[2]{\par\noindent\textsf{\textbf{#1.}} #2\par}
\tikzset{every picture/.append style={line width=1pt}}

\newtheorem{theorem}{Theorem}[chapter]
\newtheorem{proposition}[theorem]{Proposition}
\theoremstyle{definition}
\newtheorem{definition}[theorem]{Definition}
\newtheorem{example}[theorem]{Example}
\newcommand{\RR}{\mathbb{R}}
\newcommand{\CC}{\mathbb{C}}
\newcommand{\ZZ}{\mathbb{Z}}
\newcommand{\HH}{\mathbb{H}}
\newcommand{\OO}{\mathbb{O}}
\titleformat{\chapter}[display]{\normalfont\bfseries\raggedright}{\small\sffamily\chaptertitlename\ \thechapter}{0.5em}{\huge}
\titlespacing*{\chapter}{0pt}{0pt}{18pt}
\titleformat{\section}{\normalfont\Large\bfseries\raggedright}{\thesection}{0.7em}{}
\setlength{\parindent}{0pt}
\setlength{\parskip}{0.45em plus 0.1em minus 0.1em}
\setlist{itemsep=0.25em,topsep=0.4em}
\pagestyle{fancy}
\fancyhf{}
\fancyhead[L]{\small Mathematical physics / learning routes}
\fancyhead[R]{\small\thepage}
\fancyfoot[L]{\small\hyperref[learn:routes]{Routes} / \hyperref[learn:solutions]{Solutions}}
\fancyfoot[R]{\small\href{https://github.com/Oichkatzelesfrettschen/MathScienceCompendium/blob/main/docs/learning/README.md}{Companion directory}}
\raggedbottom
\hypersetup{pdftitle={From Precalculus to Mathematical Physics: A Learning Album},
 pdfauthor={Eirikr Hinngart},pdfsubject={Prerequisites, derivations, historical sources, and branching reading routes}}
\begin{document}
\frontmatter
\begin{titlepage}
\color{learnink}
\vspace*{0.12\textheight}
{\color{learnnavy}\rule{\linewidth}{1.2pt}}\par
\vspace{1em}
{\Huge\bfseries From Precalculus to\\[0.3em]Mathematical Physics\par}
\vspace{1em}
{\Large A learning album of connected mathematical journeys\par}
\vspace{1em}
{\large Derivations, historical sources, worked examples, and choices\par}
\vfill
{\large Eirikr Hinngart\par}
\medskip
A companion to \emph{Precalculus Through Problems of the Past} and
\emph{Claim, Evidence, and Falsification in a Mathematical-Physics Compendium}.
\end{titlepage}
\chapter*{What should I learn next, and why?}
\label{learn:routes}
\addcontentsline{toc}{chapter}{What should I learn next, and why?}
Begin with the question you want to answer. The independent companion,
\emph{Precalculus Through Problems of the Past}, teaches arithmetic, equations,
functions, angles, growth, and finite sums through historical problems.
Its final \emph{Choose a Method} section asks you to select an operation,
state a domain, and check units. These chapters pick up from that work.
Coordinate vectors, matrix operations, and complex arithmetic receive their
own entry lessons here.

\section*{Choose a purpose}
\textbf{Build foundations.} Open the companion's method finder and work toward
\emph{Choose a Method}. Use the entry check below to identify a useful return.
\par\textbf{Study mathematical physics.} Begin with proof, then choose a topic.
A passed entry check lets you use skills you already have. Required lessons
supply operations; helpful background supplies context.
\par\textbf{Evaluate research claims.} Start with \emph{Scope, Questions, and
Review Method} in the separate critical review. Read an overview first, then
use proof, probability, or a topic lesson to repair a specific reasoning gap.
Quantum states are a topic choice, rather than an entrance requirement for
reading the evidence method.

\section*{The companion handoff}
A rectangle has perimeter 28 and area 45. Name a side $x$, state its domain,
and find the other side. Try the operation before reading the next paragraph.
\par\textbf{Check.} The other side is $14-x$, with $0<x<14$.
The area equation gives $x^2-14x+45=(x-5)(x-9)=0$, so the sides are
5 and 9. Both measurements check. If setting up the equation is unfamiliar,
return to \emph{Equations and Balance} or \emph{Choose a Method} in the
companion. If the algebra is clear, the proof chapter asks why the argument
covers the stated domain.

\section*{Keep the books beside one another}
The album directory links separate PDFs and names each destination by book,
section, and physical PDF page. Use \emph{index.html} in the built album folder
for an offline directory, or the repository's learning guide online.
PDF viewers differ in their handling of cross-file links; the visible book and
section names remain usable with each book's outline. A browser back action
returns to the directory. The running footer supplies an online return route.
The framework is an optional workbench: choose a bounded example before
installing dependencies or interpreting a computation.

% Generated by scripts/render_learning_routes.py from library.json.

\section*{What does a symmetry preserve?}

Explain the companion transfer problems, or begin with its method finder.

\href{https://github.com/Oichkatzelesfrettschen/precalc_paper}{Companion: Precalculus}; \hyperref[learn:proof]{From convincing examples to proof}; \hyperref[learn:matrix_maps]{From balanced equations to matrix maps}; \hyperref[learn:complex_numbers]{Complex numbers: multiplication, length, and conjugation}; \hyperref[learn:linear_algebra]{Linear algebra: coordinates, maps, and modes}; \hyperref[learn:algebra]{Algebraic structure: symmetry, brackets, and roots}; \hyperref[learn:hypercomplex]{Doubling changes multiplication as well as dimension}; \href{https://github.com/Oichkatzelesfrettschen/MathScienceCompendium/blob/main/docs/learning/lessons/review_algebra.md}{Review guide: Mathematical Baseline}.

\section*{How do local rules produce changing fields?}

Explain the companion transfer problems, or begin with its method finder.

\href{https://github.com/Oichkatzelesfrettschen/precalc_paper}{Companion: Precalculus}; \hyperref[learn:proof]{From convincing examples to proof}; \hyperref[learn:matrix_maps]{From balanced equations to matrix maps}; \hyperref[learn:linear_algebra]{Linear algebra: coordinates, maps, and modes}; \hyperref[learn:calculus]{Calculus: local change and accumulated change}; \hyperref[learn:multivariable]{Several variables: local maps and conservation}; \hyperref[learn:dynamics]{Dynamics: equations, modes, and trustworthy computation}; \hyperref[learn:probability]{Probability: uncertainty, evidence, and comparison}; \hyperref[learn:numerical]{Numerical reasoning: approximation with an error budget}; \hyperref[learn:fluids]{A lattice-Boltzmann step has exact moment obligations}; \href{https://github.com/Oichkatzelesfrettschen/MathScienceCompendium/blob/main/docs/learning/lessons/review_computation.md}{Review guide: Computational Evidence Portfolio}.

\section*{How do shape and complex structure constrain a function?}

Explain the companion transfer problems, or begin with its method finder.

\href{https://github.com/Oichkatzelesfrettschen/precalc_paper}{Companion: Precalculus}; \hyperref[learn:proof]{From convincing examples to proof}; \hyperref[learn:matrix_maps]{From balanced equations to matrix maps}; \hyperref[learn:complex_numbers]{Complex numbers: multiplication, length, and conjugation}; \hyperref[learn:linear_algebra]{Linear algebra: coordinates, maps, and modes}; \hyperref[learn:calculus]{Calculus: local change and accumulated change}; \hyperref[learn:multivariable]{Several variables: local maps and conservation}; \hyperref[learn:algebra]{Algebraic structure: symmetry, brackets, and roots}; \hyperref[learn:geometry]{Geometry and analysis: distance, dimension, and complex structure}; \href{https://github.com/Oichkatzelesfrettschen/MathScienceCompendium/blob/main/docs/learning/lessons/review_geometry.md}{Review guide: Modular forms and moonshine}.

\section*{How do complex states give probabilities?}

Explain the companion transfer problems, or begin with its method finder.

\href{https://github.com/Oichkatzelesfrettschen/precalc_paper}{Companion: Precalculus}; \hyperref[learn:proof]{From convincing examples to proof}; \hyperref[learn:matrix_maps]{From balanced equations to matrix maps}; \hyperref[learn:complex_numbers]{Complex numbers: multiplication, length, and conjugation}; \hyperref[learn:linear_algebra]{Linear algebra: coordinates, maps, and modes}; \hyperref[learn:probability]{Probability: uncertainty, evidence, and comparison}; \hyperref[learn:quantum]{Worked example: probabilities from a complex state}; \href{https://github.com/Oichkatzelesfrettschen/MathScienceCompendium/blob/main/docs/learning/lessons/review_evidence.md}{Review guide: Scope, Questions, and Review Method}.

\section*{Which operation should I learn next?}

Begin with arithmetic; the companion teaches the remaining entry operations.

\href{https://github.com/Oichkatzelesfrettschen/precalc_paper}{Companion: Precalculus}; \hyperref[learn:proof]{From convincing examples to proof}.

\section*{What does the evidence establish?}

Begin with the review method as orientation. Use linked proof, probability, and topic lessons when a technical argument needs repair.

\href{https://github.com/Oichkatzelesfrettschen/MathScienceCompendium/blob/main/docs/learning/lessons/review_evidence.md}{Review guide: Scope, Questions, and Review Method}; \href{https://github.com/Oichkatzelesfrettschen/MathScienceCompendium/blob/main/docs/learning/lessons/review_physical.md}{Review guide: Audit of the Physical Claims}; \href{https://github.com/Oichkatzelesfrettschen/MathScienceCompendium/blob/main/docs/learning/lessons/review_hypotheses.md}{Review guide: Hypothesis Registry and Promotion Gates}; \href{https://github.com/Oichkatzelesfrettschen/MathScienceCompendium/blob/main/docs/learning/lessons/review_conclusions.md}{Review guide: Conclusions and Decision Record}; \href{https://github.com/Oichkatzelesfrettschen/MathScienceCompendium/blob/main/docs/learning/lessons/review_reproduction.md}{Review guide: Reproducibility and Artifact Inventory}; \href{https://github.com/Oichkatzelesfrettschen/MathScienceCompendium/blob/main/docs/learning/lessons/review_framework.md}{Review guide: Reconciled Framework Architecture}; \href{https://github.com/Oichkatzelesfrettschen/MathScienceCompendium/blob/main/docs/learning/lessons/review_critique.md}{Review guide: Falsification of the Critiques}.


\section*{Read an evidence label before accepting a claim}
\claimstatus{Definition}{A convention introduces an object and its domain.}
\claimstatus{Established result}{A proof or precise source supports a conclusion
under stated assumptions.}
\claimstatus{Computational observation}{A program checks specified inputs with
an exact or approximate method. A finite sample has a finite scope.}
\claimstatus{Conjecture or open hypothesis}{A testable statement still needs its
declared support.}
\claimstatus{Falsified within the stated scope}{A counterexample or registered
negative decision rejects the tested statement in the named regime.}
Historical cards distinguish documented sources from modern teaching examples.
The chapter bibliography and the separate evidence guide provide deeper locators.
A citation identifies a source to check; its presence alone certifies neither
access nor the historical claim.

\tableofcontents
\mainmatter
\chapter{From convincing examples to proof}
\label{learn:proof}

\begin{learningcard}{Your starting equipment}
Bring integer arithmetic, algebraic identities, inequalities, functions, and
finite sums from the precalculus book. The new skill is to specify exactly
which objects a statement covers and why the reasoning covers every one of
them. The question guiding the chapter is: when does a pattern become a theorem?
\end{learningcard}

\section{A statement needs a universe}
A proposition is a statement with a truth value once its terms and domain are
fixed. ``Every number has a square root'' changes meaning when ``number''
means a nonnegative real, an arbitrary real, or a complex number. We write
$\RR$ for the real numbers and $\ZZ$ for the integers. The symbols
$\forall$ and $\exists$ mean ``for every'' and ``there exists.'' Thus
\[
 \forall x\in\RR\;\exists y\in\RR:\quad y>x
\]
says that each real number has a larger real number. The choice $y=x+1$
proves the assertion. Reversing the quantifiers asks for one real number
larger than every real number. Any proposed $y$ fails at $x=y+1$.
The order records which choices may depend on earlier choices.

An implication $P\Rightarrow Q$ says that whenever the hypothesis $P$ holds,
the conclusion $Q$ holds. Its converse $Q\Rightarrow P$ is a separate claim.
For integers, divisibility by four implies evenness; the even integer six
refutes the converse. An equivalence $P\Leftrightarrow Q$ requires both
implications. A theorem is a proved proposition within stated assumptions;
a conjecture is a proposition offered for investigation.

\begin{learningcard}{Historical situation: a diagram and its obligations}
Euclid's \emph{Elements}, Book I, Proposition 47 proves the square relation
for a right triangle through constructions and earlier propositions. The
surviving text gives a structured chain of dependencies rather than a table
of measured triangles. Its transmission, translation, and editorial history
are separate historical questions. We use the proposition as evidence of
an organized deductive presentation, rather than as evidence that one person
invented every result in the collection. David Joyce's mathematical edition
provides the proposition, diagram, and commentary.\footnote{Euclid,
\emph{Elements}, I.47, edition and commentary by David E. Joyce, Clark
University: \url{https://mathcs.clarku.edu/~djoyce/elements/bookI/propI47.html}.}
The modern logical notation in this chapter is our teaching apparatus.
\end{learningcard}

\section{Worked example: move from a definition to a conclusion}
An integer $m$ is even when $m=2a$ for some $a\in\ZZ$, and odd when
$m=2a+1$ for some $a\in\ZZ$. Prove that the product of two odd integers
is odd. Start with arbitrary odd integers, writing $m=2a+1$, $n=2b+1$.
Then
\[
 mn=(2a+1)(2b+1)=4ab+2a+2b+1=2(2ab+a+b)+1.
\]
The expression $2ab+a+b$ is an integer because integers are closed under
addition and multiplication. The final expression therefore satisfies the
definition of oddness. The proof's decisive move is the extraction of an
integer witness, rather than the expansion alone.

A proof by contrapositive establishes $P\Rightarrow Q$ by proving
$\neg Q\Rightarrow\neg P$. For example, if $m^2$ is even, then $m$ is even:
if $m$ were odd, the product result would make $m^2$ odd. Every integer
is either even or odd, which completes the implication. A proof by
contradiction instead assumes the desired conclusion false together with
the hypotheses and derives an impossibility. These are related logical
strategies, but the argument should state which assumption it introduces.

\section{Worked example: induction is a dependency chain}
To prove a statement $P(n)$ for every integer $n\geq1$, mathematical
induction requires a base case $P(1)$ and an implication
$P(k)\Rightarrow P(k+1)$ for every integer $k\geq1$. The implication
propagates truth from the base through the positive integers.
Consider
\[
 1+2+\cdots+n=\frac{n(n+1)}2.
\]
At $n=1$, both sides equal one. Assume the formula for an arbitrary
$k\geq1$. Add the next term to obtain
\[
 \sum_{j=1}^{k+1}j=\frac{k(k+1)}2+(k+1)
 =\frac{(k+1)(k+2)}2.
\]
The right side is precisely the formula at $k+1$. The assumption covers
one unspecified preceding case; induction then supplies all cases.
Checking the first thousand cases alone would leave every later integer
outside the checked set.

The base case carries real information. The implication
``if $2^k$ is odd, then $2^{k+1}$ is odd'' holds for $k\geq1$ because
its hypothesis always fails, yet $2^n$ is even for every positive $n$.
A chain of implications needs a true starting point.

\section{A counterexample has a small job}
A universal claim fails when one admissible object violates its conclusion.
For ``every prime is odd,'' two is sufficient: it is prime and even.
A counterexample must satisfy the hypotheses. The number one fails as a
counterexample because a prime is an integer greater than one whose only
positive divisors are one and itself.

An existential statement needs one witness. A universal statement over a
finite, completely enumerated domain can be proved by exact exhaustive
calculation. The same calculation over a sample from an infinite domain
establishes only the sampled cases. Floating-point computations introduce
a further boundary: a small residual measures agreement to a tolerance,
whereas exact equality requires an exact argument or exact arithmetic.

\begin{learningcard}{Choose a useful next question}
Before calculating, write the domain, quantifiers, and assumptions. After
calculating, identify the witness or implication that makes the result
apply to its claimed domain. Continue to Chapter~\ref{learn:linear_algebra}
for structural proofs or Chapter~\ref{learn:calculus} for quantified limits.
\end{learningcard}

\section{Exercises with full solutions}
\paragraph{1. Negate a quantified statement.}
Negate: ``For every real $x$ there exists a real $y$ such that $y^2=x$.''
Determine the truth of the original and its negation.

\paragraph{Solution.}
The negation is ``There exists a real $x$ such that for every real $y$,
$y^2\ne x$.'' Choose $x=-1$. Every real square is nonnegative, so every
real $y$ satisfies $y^2\ne-1$. The negation is true and the original is
false. The witness refutes a universal statement; checking positive
values of $x$ would leave the negative part of the domain unresolved.

\paragraph{2. Prove an inequality by induction.}
Prove $2^n\geq n+1$ for every integer $n\geq0$.

\paragraph{Solution.}
At zero, $2^0=1=0+1$. Assuming $2^k\geq k+1$ for an arbitrary
$k\geq0$, multiplication by the positive number two gives
$2^{k+1}\geq2k+2$. Since $2k+2-(k+2)=k\geq0$, we obtain
$2^{k+1}\geq k+2$. Induction completes the proof, including the stated
zero endpoint.

\paragraph{3. Separate an implication and its converse.}
For a real number $x$, analyze ``$x>2$ implies $x^2>4$'' and its converse.

\paragraph{Solution.}
If $x>2$, then both $x-2$ and $x+2$ are positive, so
$x^2-4=(x-2)(x+2)>0$. The converse fails at $x=-3$, whose square is
nine. A corrected equivalence is $x^2>4$ if and only if $|x|>2$:
the factors show the quadratic is positive on $(-\infty,-2)$ and
$(2,\infty)$, exactly the domain described by the absolute-value inequality.

\chapter{From balanced equations to matrix maps}
\label{learn:matrix_maps}
\begin{learningcard}{Bring forward and choose your next operation}
\textsf{\textbf{Question.}} How can two accounts become one operation?\par
\textbf{Required lessons.} \hyperref[learn:proof]{From convincing examples to proof}\par
\textbf{Helpful background.} Spatial arrows help interpret a vector; units determine whether a norm is a physical length.\par
\textbf{Learn to.} Compute matrix actions, compositions, and real dot products with compatible dimensions.\par
\textbf{Entry check.} Solve x+y=7 and 2x+y=11.\par
\hyperref[solution:entry:matrix_maps]{Check your answer}; repair the operation in \href{https://github.com/Oichkatzelesfrettschen/precalc_paper}{Companion: Precalculus}.
\end{learningcard}


\section{How can two accounts become one operation?}
The companion's balanced-equation lesson recovers two yields from two
accounts. Keep that question, and change the representation. Suppose two
bundles of one grade and one of another yield eleven measures, while one
of the first and two of the second yield ten. These are modern teaching
accounts. Each grade has a fixed yield in measures per bundle.
\lessoncue{Notice}{The same two yields occur in both accounts.}
Write their numerical values in an ordered column
\[
 v=\begin{pmatrix}x\\y\end{pmatrix},\qquad
 A=\begin{pmatrix}2&1\\1&2\end{pmatrix},\qquad
 Av=\begin{pmatrix}2x+y\\x+2y\end{pmatrix}.
\]
The entries of $A$ count bundles; multiplying bundle counts by yields
produces measures. An ordered column of real numbers is a coordinate
vector. Order carries meaning: exchanging $x$ and $y$ exchanges the grades.
A matrix is a rectangular array; here its two rows describe the accounts.

\section{Combine corresponding coordinates}
For columns of equal length, add corresponding entries and multiply every
entry by a scalar:
\[
 \begin{pmatrix}4\\3\end{pmatrix}+\begin{pmatrix}1\\-1\end{pmatrix}
 =\begin{pmatrix}5\\2\end{pmatrix},\qquad
 2\begin{pmatrix}4\\3\end{pmatrix}=\begin{pmatrix}8\\6\end{pmatrix}.
\]
Subtraction adds the negative column. The zero column leaves every entry
unchanged under addition. A linear combination such as $2v-w$ means scale
$v$, then subtract corresponding coordinates of $w$. In a physical model,
corresponding quantities must have compatible units before addition.

\section{Multiply a row by a column}
For a matrix with $m$ rows and $n$ columns, an input column must have $n$
entries. Multiply matching entries of a row and the input, then add; repeat
for each row. The output has $m$ entries. In symbols,
$(Av)_r=\sum_{s=1}^n A_{rs}v_s$: the summation sign asks you to add the
products as $s$ runs from one to $n$.
\lessoncue{Worked example}{Check the recovered yields in both accounts.}
Substitute $x=4,y=3$:
\[
 A\begin{pmatrix}4\\3\end{pmatrix}
 =\begin{pmatrix}2\cdot4+1\cdot3\\1\cdot4+2\cdot3\end{pmatrix}
 =\begin{pmatrix}11\\10\end{pmatrix}.
\]
Solving $Av=(11,10)^T$ asks which input gives that output. The superscript
$T$ transposes an array: rows become columns. Thus $(11,10)^T$ is a column.
For a rectangular example,
\[
 \begin{pmatrix}1&2&3\\4&5&6\end{pmatrix}^{T}
 =\begin{pmatrix}1&4\\2&5\\3&6\end{pmatrix}.
\]
Transpose changes positions, while applying a matrix computes new values.

\begin{figure}[htbp]
\centering
\begin{tikzpicture}[>=Stealth,font=\small]
\node[draw,align=center,text width=2.8cm] (input) at (0,0)
 {$v=(4,3)^T$\\yields\\measures per bundle};
\node[draw,align=center,text width=2.8cm] (accounts) at (4.5,0)
 {$Av=(11,10)^T$\\account amounts\\measures};
\node[draw,align=center,text width=2.3cm] (difference) at (8.8,0)
 {$BAv=1$\\difference\\measure};
\draw[->] (input)--node[above] {$A$} (accounts);
\draw[->] (accounts)--node[above] {$B$} (difference);
\draw[->] (input.south) -- ++(0,-0.7) -|
 node[pos=0.25,below] {$BA$: subtract rows, then multiply} (difference.south);
\end{tikzpicture}

\caption{Read each arrow as an operation. The first produces two account
amounts; the second takes their difference. The composite row computes
that difference directly from the yields. Units travel through the arrows.}
\end{figure}

\section{Compose operations in their application order}
Take the first account minus the second with the row $B=(1,-1)$.
Then $B(Av)=11-10=1$ measure. Distribute the subtraction before inserting
numbers: $(2x+y)-(x+2y)=x-y$. Therefore
\[
 BA=\begin{pmatrix}1&-1\end{pmatrix}
     \begin{pmatrix}2&1\\1&2\end{pmatrix}
   =\begin{pmatrix}1&-1\end{pmatrix},\qquad B(Av)=(BA)v.
\]
The rightmost matrix acts first. In general, if $A$ has size $m\times n$
and $B$ has size $p\times m$, then $BA$ has size $p\times n$.
Its entry in row $r$, column $s$ is the row $r$ of $B$ multiplied by
column $s$ of $A$. The inner sizes must agree. Here $AB$ is undefined:
$A$ expects two input coordinates, while $B$ produces one.

Even when both products exist, their order can matter. Let
\[
 C=\begin{pmatrix}1&1\\0&1\end{pmatrix},\quad
 D=\begin{pmatrix}2&0\\0&1\end{pmatrix}.
\]
Row-column multiplication gives $CD=\begin{pmatrix}2&1\\0&1\end{pmatrix}$
and $DC=\begin{pmatrix}2&2\\0&1\end{pmatrix}$.
Doubling the first coordinate before adding the second differs from doubling
after the addition. The identity matrix $I$, with ones on its diagonal and
zeros elsewhere, satisfies $Iv=v$.

\section{Measure length only after choosing a meaning}
For real columns of equal length define the Euclidean dot product and norm:
\[
 u\mathbin{\cdot}v=u^Tv=\sum_{s=1}^n u_sv_s,\qquad
 \|v\|=\sqrt{v^Tv}=\sqrt{\sum_{s=1}^n v_s^2}.
\]
For displacement coordinates on perpendicular axes in the same length unit,
the Pythagorean theorem identifies this norm as length. For unrelated
quantities, choose scales and a physical interpretation before calling a
coordinate norm a measured distance.
\lessoncue{Worked example}{Find a unit direction and a perpendicular one.}
For $v=(3,4)^T$, the squared norm is $9+16=25$, so $\|v\|=5$.
Divide each coordinate by five to obtain $q=(3/5,4/5)^T$ with norm one.
For $w=(-4,3)^T$, $v^Tw=-12+12=0$: the directions are perpendicular.
A zero vector has norm zero; division by its norm is undefined.

\section{Try and check}
\paragraph{1. Transfer the account rule.}
For $M=\begin{pmatrix}3&1\\1&2\end{pmatrix}$ and $v=(2,5)^T$,
find $Mv$ and check the difference of its accounts directly.
\paragraph{Solution.}
Multiply each row: $Mv=(3\cdot2+5,2+2\cdot5)^T=(11,12)^T$.
Subtract the outputs to get $-1$. Subtract the rows first to obtain
$(2,-1)$, then compute $2\cdot2-5=-1$. Both orders give the same difference.

\paragraph{2. Check dimensions and order.}
Using $C,D$ above and $u=(1,2)^T$, calculate $C(Du)$ and $D(Cu)$.
\paragraph{Solution.}
First $Du=(2,2)^T$, then $C(Du)=(4,2)^T$.
First $Cu=(3,2)^T$, then $D(Cu)=(6,2)^T$. The intermediate columns
have two entries, so both operations are defined, but their outputs differ.

\paragraph{3. Normalize and test perpendicularity.}
Find the norm of $u=(5,12)^T$ and a perpendicular unit vector.
\paragraph{Solution.}
The squared norm is $25+144=169$, giving norm thirteen.
The vector $w=(-12/13,5/13)^T$ has squared norm
$(144+25)/169=1$ and $u^Tw=(-60+60)/13=0$.

\section*{Ready for the next question?}
\label{readiness:matrix_maps}
Apply the matrix with rows (3,1) and (1,2) to (2,5), then subtract the second output from the first.
\par\hyperref[solution:ready:matrix_maps]{Read the worked readiness solution.} Continue when you can explain the operations and assumptions. Otherwise return to the method and its solved exercises.

\lessonreference
{Match row and column sizes; multiply corresponding entries and add.
Apply $A$ before $B$ in $BA$. For real coordinates, $\|v\|^2=v^Tv$.}
{Coordinate order and units carry meaning. Transpose swaps rows and columns.
Matrix multiplication can depend on order; normalization requires $v\ne0$.}
{Chapter~\ref{learn:linear_algebra} uses these operations to find independent
coordinates and eigenvectors. Chapter~\ref{learn:complex_numbers} extends
length calculations to complex coordinates.}


\chapter{Complex numbers: multiplication, length, and conjugation}
\label{learn:complex_numbers}
\begin{learningcard}{Bring forward and choose your next operation}
\textsf{\textbf{Question.}} How does multiplication change a complex number?\par
\textbf{Required lessons.} \hyperref[learn:proof]{From convincing examples to proof}\par
\textbf{Helpful background.} A coordinate-plane sketch helps interpret multiplication.\par
\textbf{Learn to.} Compute conjugates, lengths, division, and normalized complex columns.\par
\textbf{Entry check.} Expand (a+b)(c+d).\par
\hyperref[solution:entry:complex_numbers]{Check your answer}; repair the operation in \href{https://github.com/Oichkatzelesfrettschen/precalc_paper}{Companion: Precalculus}.
\end{learningcard}


\section{How can arithmetic beyond real roots retain a positive length?}
The companion's completing-the-square lesson introduces $i^2=-1$ to solve
an equation such as $x^2+4=0$. Carry that arithmetic forward: a complex
number $z=a+bi$ has real coordinates $a,b$. Equality means equality of
both coordinates. A real number embeds as $a+0i$.
Our question now concerns measurement. Squaring $i$ gives $-1$, yet a
squared length should be nonnegative. Which product measures length?

\section{Distribute, then replace the square of \texorpdfstring{$i$}{i}}
Add and subtract corresponding real and imaginary parts. For multiplication,
use distribution before collecting terms:
\[
 (a+bi)(c+di)=ac+adi+bci+bdi^2=(ac-bd)+(ad+bc)i.
\]
\lessoncue{Worked example}{Separate the arithmetic operations.}
For $z=2+i$ and $w=1-3i$, addition gives $z+w=3-2i$.
Multiplication gives
\[
 zw=2-6i+i-3i^2=5-5i.
\]
The second equality replaces $i^2$ by $-1$ and combines imaginary terms.
Multiplying a general number by $i$ gives $i(a+bi)=-b+ai$: its plane
coordinates change from $(a,b)$ to $(-b,a)$, a quarter-turn counterclockwise.

\begin{figure}[htbp]
\centering
\begin{tikzpicture}[x=1.25cm,y=1.25cm,>=Stealth,font=\small]
\draw[->] (-1.8,0)--(3,0) node[right] {real};
\draw[->] (0,-1.6)--(0,2.7) node[above] {imaginary};
\draw[->,thick] (0,0)--(2,1) node[right] {$z=2+i$};
\draw[->,thick,dashed] (0,0)--(-1,2) node[left] {$iz=-1+2i$};
\draw[->,thick,dotted] (0,0)--(2,-1) node[right] {$\overline z=2-i$};
\draw[->] (0.55,0.275) arc[start angle=26.565,end angle=116.565,radius=0.615];
\node at (-0.05,0.85) {$\times i$};
\foreach \position in {-1,1,2} {
\draw (\position,0.05)--(\position,-0.05) node[below] {$\position$};
\draw (0.05,\position)--(-0.05,\position);
}
\end{tikzpicture}

\caption{The solid arrow represents $z=2+i$. Multiplication by $i$ gives
the dashed arrow $iz=-1+2i$, a quarter-turn of equal length. Reflection
across the real axis gives the dotted arrow $\overline z=2-i$. Line styles
and endpoint labels distinguish the three operations in grayscale.}
\end{figure}

\section{Conjugate to measure length and divide}
The conjugate of $z=a+bi$ is $\overline z=a-bi$. Multiply the pair:
\[
 \overline z z=(a-bi)(a+bi)=a^2+b^2.
\]
The cross terms cancel. Define the modulus $|z|=\sqrt{a^2+b^2}$,
the distance from the origin in the complex plane. Therefore
$|z|^2=\overline z z$, and $|z|=0$ exactly when $z=0$.
For $z=2+i$, $|z|=\sqrt5$, while $z^2=3+4i$ is a different quantity.

\lessoncue{Worked example}{Make a denominator real before dividing.}
For $w\ne0$, multiply numerator and denominator by $\overline w$.
The denominator becomes the positive real number $|w|^2$:
\[
 \frac{2+i}{1-3i}
 =\frac{(2+i)(1+3i)}{(1-3i)(1+3i)}
 =\frac{-1+7i}{10}=-\frac1{10}+\frac7{10}i.
\]
Check by multiplying the answer by $1-3i$:
$(-1+7i)(1-3i)/10=(20+10i)/10=2+i$.
Distribution also verifies $\overline{zw}=\overline z\,\overline w$.
Consequently $|zw|^2=|z|^2|w|^2$, so $|zw|=|z||w|$ because
moduli are nonnegative. Multiplication by $i$ preserves length since $|i|=1$.

\section{Extend the dot product with conjugation}
A complex coordinate vector is a column whose entries are complex numbers.
Add columns component by component and scale each entry as before. For
equal-length columns define
\[
 \langle u,v\rangle=\sum_{s=1}^n\overline{u_s}v_s,
 \qquad \|v\|=\sqrt{\langle v,v\rangle}
 =\sqrt{\sum_{s=1}^n |v_s|^2}.
\]
We conjugate the first argument, the convention used in the quantum bridge.
Thus $\langle u,cv\rangle=c\langle u,v\rangle$, while
$\langle cu,v\rangle=\overline c\langle u,v\rangle$.
Exchanging the arguments conjugates the result:
$\langle v,u\rangle=\overline{\langle u,v\rangle}$.
The conjugate transpose $u^*$ means transpose the column and conjugate
its entries, so $\langle u,v\rangle=u^*v$; $u^\dagger$ is another common
notation. For real columns the formula reduces to the ordinary dot product.

\lessoncue{Worked example}{Test two complex directions.}
Take $u=(1,i)^T$ and $v=(i,1)^T$. Then
\[
 \langle u,v\rangle=1\cdot i+(-i)\cdot1=0,\qquad
 \|u\|^2=1+(-i)i=2.
\]
The columns are orthogonal, meaning their inner product is zero.
Dividing either column by $\sqrt2$ gives a unit vector. Omitting conjugation
would give $u^Tu=1+i^2=0$ for a nonzero column, so that expression fails
as a squared norm. In quantum mechanics the normalized column
$u/\sqrt2$ has squared component moduli $1/2,1/2$. Their interpretation as
measurement probabilities requires the state and measurement assumptions
introduced in Chapter~\ref{learn:quantum}.

\section{Try and check}
\paragraph{1. Multiply and check the length rule.}
For $z=3-2i$ and $w=-1+i$, find $zw$ and verify $|zw|^2=|z|^2|w|^2$.
\paragraph{Solution.}
Distribution gives $zw=-3+3i+2i-2i^2=-1+5i$.
Its squared modulus is $1+25=26$. The separate squared moduli are
$9+4=13$ and $1+1=2$, whose product is twenty-six.

\paragraph{2. Divide and reconstruct the numerator.}
Compute $(3+i)/(1+i)$.
\paragraph{Solution.}
Multiply both numerator and denominator by $1-i$:
$(3+i)(1-i)/2=(4-2i)/2=2-i$.
Multiplication checks $(2-i)(1+i)=3+i$.

\paragraph{3. Normalize a complex column.}
For $v=(1+i,2)^T$, find a unit vector and verify its squared norm.
\paragraph{Solution.}
The squared norm is $|1+i|^2+|2|^2=2+4=6$.
Set $q=v/\sqrt6$. Its squared norm is $2/6+4/6=1$.
Each component uses its modulus squared; squaring the complex component
itself would give $(1+i)^2=2i$ and lose the length interpretation.

\section*{Ready for the next question?}
\label{readiness:complex_numbers}
Find the conjugate and squared modulus of 3+4i, then its reciprocal.
\par\hyperref[solution:ready:complex_numbers]{Read the worked readiness solution.} Continue when you can explain the operations and assumptions. Otherwise return to the method and its solved exercises.

\lessonreference
{Distribute products using $i^2=-1$; use $\overline z z=|z|^2$ for
length. Divide by $w\ne0$ by multiplying by $\overline w/\overline w$.}
{The complex inner product conjugates its first argument here.
Normalization divides a nonzero vector by its positive real norm.}
{Chapter~\ref{learn:quantum} uses complex amplitudes and inner products;
Chapter~\ref{learn:geometry} uses lengths and products as geometric data.}


\chapter{Linear algebra: coordinates, maps, and modes}
\label{learn:linear_algebra}

\begin{learningcard}{Prerequisites and destination}
Bring precalculus vectors, dot products, matrix multiplication, and systems
of equations; use Chapter~\ref{learn:proof} for quantified statements.
We ask how a coupled problem can be expressed in coordinates that expose
independent directions. Eigenvectors will later become the modes of
Chapter~\ref{learn:dynamics}.
\end{learningcard}

\section{Vectors extend beyond arrows}
A real vector space is a set $V$ with addition and multiplication by real
scalars satisfying the following rules: addition is associative and
commutative, a zero vector and additive inverses exist, scalar multiplication
is associative with scalar products, one acts as the identity, and the two
distributive laws hold. The operations stay inside $V$.
The familiar space $\RR^n$ satisfies these rules component by component.
Polynomials of degree at most two form another example: adding two such
polynomials and scaling one preserves that degree bound.

A linear combination of vectors $v_1,\ldots,v_k$ is a sum
$\sum_{j=1}^k a_jv_j$ with real coefficients. Their span is the set of all
such combinations. They are linearly independent when
$\sum a_jv_j=0$ forces every $a_j=0$. A basis is an independent spanning
list. Every vector has unique coefficients in a basis: subtracting two
representations would give a vanishing linear combination, so independence
forces the coefficients to agree. The number of vectors in a basis of a
finite-dimensional space is its dimension.

For degree-at-most-two polynomials, $(1,x,x^2)$ is a basis. Coordinates
$(3,-2,5)$ represent $3-2x+5x^2$. Calling a polynomial a vector identifies
its algebraic operations; magnitude and spatial direction require additional
structure and interpretation.

\section{Worked example: a map becomes a matrix}
A map $T:V\to W$ is linear if
$T(au+bv)=aT(u)+bT(v)$ for all vectors and real scalars. Define
$T(a+bx+cx^2)=b+2cx$ on degree-at-most-two polynomials. This coefficient
rule defines a linear map using only algebra. Chapter~\ref{learn:calculus}
identifies the same rule as polynomial differentiation, $T(p)=p'$. Regard
the output as another polynomial of degree at most two. Since
$T(1)=0$, $T(x)=1$, and $T(x^2)=2x$, its matrix in the basis above is
\[
 A=\begin{pmatrix}0&1&0\\0&0&2\\0&0&0\end{pmatrix}.
\]
Each column is the coordinate vector of a basis image. Multiplication gives
$A(3,-2,5)^T=(-2,10,0)^T$, representing $-2+10x$ as required.
The kernel is the set of vectors mapped to zero, here the constant
polynomials. The image is the set of outputs, here polynomials of degree
at most one. Their dimensions are one and two.

The rank--nullity theorem states
$\dim\ker T+\dim\operatorname{im}T=\dim V$ for finite-dimensional $V$.
To see the mechanism, take a basis of the kernel and extend it to a basis
of $V$. The images of the added basis vectors span the image. Any linear
relation among those images would put a combination of the added vectors
inside the kernel, contradicting independence of the extended basis.
Thus the added images form a basis of the image.

\begin{learningcard}{Historical situation: an algebra of transformations}
Arthur Cayley's 1858 ``A Memoir on the Theory of Matrices'' treats matrices
as objects with operations and develops their algebra. The memoir is a
specific document in the nineteenth-century development of the subject;
methods for solving linear systems have much older histories. The Royal
Society publication context shows mathematical communication through a
learned institution, rather than a solitary discovery detached from readers
and prior techniques.\footnote{Arthur Cayley, ``A Memoir on the Theory of
Matrices,'' \emph{Philosophical Transactions of the Royal Society of London}
148 (1858), pp.~17--37, especially the opening definitions:
\url{https://doi.org/10.1098/rstl.1858.0002}.}
The polynomial example is a modern instructional construction.
\end{learningcard}

\section{Worked example: independent stretching directions}
An eigenvector of a square matrix $A$ is a nonzero vector $v$ satisfying
$Av=\lambda v$; its eigenvalue $\lambda$ is the stretching factor,
including a reversal when negative. For
\[
 A=\begin{pmatrix}2&1\\1&2\end{pmatrix},\qquad
 v_+=(1,1)^T,\qquad v_-=(1,-1)^T,
\]
direct multiplication gives $Av_+=3v_+$ and $Av_-=v_-$. These
independent vectors form a basis. Since
$(4,2)^T=3v_++v_-$, repeated application gives
$A^m(4,2)^T=3\cdot3^m v_++v_-$ for integers $m\geq0$.
A coupled recurrence becomes two scalar recurrences.

Normalize the vectors to obtain $q_+=v_+/\sqrt2$ and
$q_-=v_-/\sqrt2$. The matrix $Q$ with those columns obeys
$Q^TQ=I$, where $I$ is the identity matrix. Such a matrix is orthogonal;
its columns form an orthonormal basis, meaning unit lengths and pairwise
zero dot products. In these coordinates,
$Q^TAQ=\operatorname{diag}(3,1)$ and
$A=Q\operatorname{diag}(3,1)Q^T$.

\begin{figure}[htbp]
  \centering
  \begin{tikzpicture}[x=0.85cm,y=0.85cm,>=Stealth,font=\small,line width=0.8pt]
  \begin{scope}
    \draw[->] (-0.4,0)--(3.6,0) node[right] {$x$};
    \draw[->] (0,-1.5)--(0,3.6) node[above] {$y$};
    \node[above] at (1.7,4.1) {Input directions};
    \draw[->,line width=1.2pt] (0,0)--(1,1)
      node[above right] {$v_+=(1,1)^T$};
    \draw[->,dashed,line width=1.2pt] (0,0)--(1,-1)
      node[below right] {$v_-=(1,-1)^T$};
    \node[below left] at (0,0) {$0$};
    \foreach \value in {1,2,3} {
      \draw (\value,0.07)--(\value,-0.07);
      \draw (0.07,\value)--(-0.07,\value);
    }
  \end{scope}
  \begin{scope}[xshift=6.2cm]
    \draw[->] (-0.4,0)--(3.6,0) node[right] {$x$};
    \draw[->] (0,-1.5)--(0,3.6) node[above] {$y$};
    \node[above] at (1.7,4.1) {After applying $A$};
    \draw[->,line width=1.2pt] (0,0)--(3,3)
      node[above left] {$Av_+=3v_+$};
    \draw[->,dashed,line width=1.2pt] (0,0)--(1,-1)
      node[below right] {$Av_-=v_-$};
    \node[below left] at (0,0) {$0$};
    \foreach \value in {1,2,3} {
      \draw (\value,0.07)--(\value,-0.07) node[below] {$\value$};
      \draw (0.07,\value)--(-0.07,\value);
    }
  \end{scope}
  \node at (5.3,2) {$\displaystyle A=\begin{pmatrix}2&1\\1&2\end{pmatrix}$};
\end{tikzpicture}

  \caption{Read the input arrows at left, then the corresponding arrows at right on the same coordinate scale. The solid eigenvector grows by a factor of three; the dashed eigenvector keeps its length and direction. A general input splits into these two directions.}
  \label{fig:learn:linear_map}
\end{figure}

\section{Why symmetry guarantees enough eigenvectors}
The real symmetric spectral theorem states that every real matrix satisfying
$A^T=A$ has an orthonormal basis of real eigenvectors. Here is a proof
outline, with its analysis input made explicit. The continuous function
$q(v)=v^TAv$ attains a maximum on the unit sphere because that sphere is
compact in finite-dimensional Euclidean space. Let $u$ be a maximizer.
For any unit $w$ perpendicular to $u$, the curve
$u\cos t+w\sin t$ stays on the sphere. Differentiating $q$ along the
curve at $t=0$ gives $2w^TAu=0$. Consequently $Au$ has zero component
perpendicular to $u$, so $Au=\lambda u$.

For $w\perp u$, symmetry gives $u^TAw=(Au)^Tw=\lambda u^Tw=0$.
Thus the orthogonal complement of $u$ is preserved by $A$, and the
restriction there is again symmetric. Induction on dimension supplies an
orthonormal eigenbasis of that complement, which together with $u$ proves
the theorem. The compactness and derivative facts belong to real analysis;
the remainder is linear algebra.

A general real matrix need not have real eigenvectors: rotation by ninety
degrees sends every nonzero vector to a perpendicular vector. Even over
$\CC$, the matrix
\[J=\begin{pmatrix}1&1\\0&1\end{pmatrix}\]
has only a one-dimensional eigenspace and lacks an eigenbasis. These examples
explain why the theorem explicitly requires real symmetry.

\section{Exercises with full solutions}
\paragraph{1. Test a candidate vector space.}
Is $S=\{(x,y)\in\RR^2:x+y=1\}$ a subspace with ordinary operations?

\paragraph{Solution.}
A subspace must contain the zero vector, but $0+0\ne1$. Also $(1,0)$
belongs to $S$ while twice that vector does not. Therefore $S$ is an affine
line rather than a vector subspace. The parallel set $x+y=0$ is a subspace:
it equals the span of $(1,-1)$ and has dimension one.

\paragraph{2. Find a projection matrix and its modes.}
Find the orthogonal projection onto the line spanned by $(1,1)$.

\paragraph{Solution.}
A unit direction is $q=(1,1)^T/\sqrt2$. The projected vector is
$q(q^Tv)$, so
\[P=qq^T=\frac12\begin{pmatrix}1&1\\1&1\end{pmatrix}.\]
Multiplication gives $P(1,1)^T=(1,1)^T$ and $P(1,-1)^T=0$.
The eigenvalues are one along the line and zero perpendicular to it.
Applying the projection again preserves each component, so $P^2=P$.

\paragraph{3. Check an energy claim.}
For the symmetric example above, show $v^TAv>0$ for every nonzero real $v$.

\paragraph{Solution.}
Write $v=aq_++bq_-$. Orthonormality and the eigenvalue equations give
$v^TAv=3a^2+b^2$. A nonzero vector has at least one nonzero coefficient,
so the expression is positive. This property is called positive definiteness.
For a real symmetric matrix it holds exactly when every eigenvalue is
positive, since the same orthonormal expansion works in every dimension.

\chapter{Calculus: local change and accumulated change}
\label{learn:calculus}

\begin{learningcard}{Prerequisites and motivating problem}
Bring functions, inequalities, finite geometric sums, and the informal limits
from precalculus; Chapter~\ref{learn:proof} supplies quantifiers. A speedometer
reports local motion, while an odometer reports accumulated distance. What
mathematics connects a rate measured near one instant to a total over an
interval? Signed displacement, rather than total distance, will be our first
quantity; total distance integrates the magnitude of velocity.
\end{learningcard}

\section{A limit is a controlled error statement}
For a function defined near $a$, except possibly at $a$, the assertion
$\lim_{x\to a}f(x)=L$ means: for every $\epsilon>0$ there exists
$\delta>0$ such that $0<|x-a|<\delta$ implies $|f(x)-L|<\epsilon$.
Inputs are restricted to the function's domain, and $a$ must be an
accumulation point of that domain. The requested output tolerance comes
first; our input tolerance may depend on it. A function is continuous at
$a$ when its limit there equals $f(a)$.

For $f(x)=3x+1$ at $a=2$, the expected limit is seven. Given
$\epsilon>0$, choose $\delta=\epsilon/3$. Then
$|f(x)-7|=3|x-2|<3\delta=\epsilon$. The proof produces a rule
for every requested tolerance, rather than a finite list of nearby values.
For a sequence, $a_n\to L$ means that for every $\epsilon>0$ there is
an integer $N$ such that $n\geq N$ implies $|a_n-L|<\epsilon$.
The input control is now a sufficiently late index.

\section{Worked example: derive a rate before memorizing a rule}
The derivative is
\[
 f'(a)=\lim_{h\to0}\frac{f(a+h)-f(a)}h
\]
when that finite two-sided limit exists. For $f(x)=x^2$, expand first:
\[
 \frac{(a+h)^2-a^2}{h}=2a+h\quad(h\ne0).
\]
Taking the limit gives $f'(a)=2a$. The corresponding tangent line is
$y=a^2+2a(x-a)$. Its error is exactly $(x-a)^2$, so the error divided
by $|x-a|$ tends to zero. Differentiability says precisely that a linear
approximation has an error smaller than first order.

Using the derivative definition, adding and subtracting
$f(a+h)g(a)$ inside the difference quotient proves the product rule
$(fg)'=f'g+fg'$. The limit uses continuity of differentiable functions.
The chain rule is $(f\circ g)'(a)=f'(g(a))g'(a)$ when both relevant
derivatives exist. These rules preserve assumptions: $|x|$ is continuous
at zero but its left and right slopes are $-1$ and $1$, so its derivative
there fails to exist.

The trigonometric rules used later are $(\sin x)'=\cos x$ and
$(\cos x)'=-\sin x$, with angles in radians. The unit-circle area
comparison gives $\sin h<h<\tan h$ for $0<h<\pi/2$, hence
$\sin h/h\to1$ by squeezing; symmetry handles negative $h$.
Also $(\cos h-1)/h=-\sin^2h/[h(1+\cos h)]\to0$.
Substitution into the angle-addition formulas proves the two derivative
rules. After the fundamental theorem below, one may define
$\log x=\int_1^x dt/t$ for $x>0$. Its derivative is $1/x$;
the derivative rule for an inverse function then gives
$(e^x)'=e^x$ for its inverse exponential. These identities support the
rates and Taylor examples used throughout the album.

\begin{learningcard}{Historical situation: motion as a mathematical problem}
In his preface to the first edition of the \emph{Principia}, Newton
connects mathematical principles with the investigation of forces from
motions and motions from forces. Book I develops mathematical propositions
about motion through a geometrical presentation. The published text supports
that methodological description; it does not make our modern epsilon--delta
notation Newton's own.\footnote{Isaac Newton, \emph{Philosophiae Naturalis
Principia Mathematica} (1687), author's preface and Book I; Cambridge
University Library digitization of Newton's copy:
\url{https://cudl.lib.cam.ac.uk/view/PR-ADV-B-00039-00001/1}.}
A history of calculus also requires Leibniz and earlier traditions; this
single source establishes the stated Newtonian context only.
\end{learningcard}

\section{Worked example: accumulated change as a limit of sums}
For a bounded function on $[a,b]$, divide the interval into pieces with
endpoints $a=x_0<\cdots<x_N=b$. Choose a sample $\xi_j$ in each piece.
The Riemann sum is $\sum_j f(\xi_j)(x_j-x_{j-1})$. The Riemann integral
is the common limit of these sums as the maximum piece width tends to zero,
independently of the partitions and samples. Continuous functions on a closed
bounded interval have such an integral. A signed integral allows negative
contributions below the horizontal axis.

For $f(x)=x$ on $[0,1]$, equal pieces and right endpoints give
\[
 \frac1N\sum_{j=1}^N\frac jN
 =\frac{N(N+1)}{2N^2}=\frac12+\frac1{2N}\longrightarrow\frac12.
\]
Continuity guarantees sample independence. More directly, the left sum is
$1/2-1/(2N)$, and monotonicity traps any sampled sum between those two
values. The limit reproduces the triangle's area while exposing the error.

\section{Why accumulation and differentiation meet}
Let $f$ be continuous on $[a,b]$ and define
$F(x)=\int_a^x f(t)\,dt$. For an interior point $x$ and small nonzero $h$,
\[
 \frac{F(x+h)-F(x)}h-f(x)
 =\frac1h\int_x^{x+h}\bigl(f(t)-f(x)\bigr)\,dt.
\]
The magnitude is at most the largest value of $|f(t)-f(x)|$ between
$x$ and $x+h$, which tends to zero by continuity. Thus $F'(x)=f(x)$.
If $G$ is any antiderivative of $f$, the derivative of $G-F$ vanishes.
The mean value theorem makes $G-F$ constant, giving the fundamental theorem
of calculus:
\[
 \int_a^b f(x)\,dx=G(b)-G(a).
\]
The mean value theorem used here states that a continuous function on a
closed interval, differentiable in its interior, has a point whose derivative
equals the interval's secant slope.

If velocity is $v(t)=3t^2$ in metres per second with time measured in seconds,
the displacement from $t=0$ to $t=2$ is $[t^3]_0^2=8$ metres.
The units explain why multiplying a rate by a time interval is necessary.

\section{Taylor polynomials need remainder control}
For $f$ with $n+1$ continuous derivatives on the interval joining $a$ and $x$,
Taylor's theorem gives
\[
 f(x)=\sum_{k=0}^n\frac{f^{(k)}(a)}{k!}(x-a)^k+R_n(x),\qquad
 |R_n(x)|\leq\frac{M|x-a|^{n+1}}{(n+1)!},
\]
where $M$ bounds $|f^{(n+1)}|$ on that interval. Here $k!$ is the
product of the positive integers through $k$, with $0!=1$.
Repeated integration of the fundamental theorem gives an integral remainder
$R_n(x)=\frac1{n!}\int_a^x f^{(n+1)}(t)(x-t)^n\,dt$;
its absolute-value bound yields the displayed estimate.

At $a=0$, the degree-three Taylor polynomial of sine is $x-x^3/6$.
Since the fourth derivative is sine and has magnitude at most one, at
$x=0.1$ radians the error is at most $0.1^4/24$.
For equality with an infinite Taylor series we need $R_n(x)\to0$.
Infinitely many derivatives alone are insufficient. The function equal to
$e^{-1/x^2}$ for $x\ne0$ and zero at zero is smooth with every derivative
zero at zero, yet is positive elsewhere. Its derivative expressions are
finite sums of inverse powers of $x$ times $e^{-1/x^2}$; exponential decay
beats every such power as $x\to0$. This proves the asserted zero derivatives
inductively. A function equal to its Taylor series locally is called analytic.

\section{Exercises with full solutions}
\paragraph{1. Supply a tolerance rule.}
Prove $\lim_{x\to2}x^2=4$ using the definition.

\paragraph{Solution.}
Factor $|x^2-4|=|x-2||x+2|$. Restrict $|x-2|<1$, which implies
$1<x<3$ and $|x+2|<5$. Given $\epsilon>0$, choose
$\delta=\min(1,\epsilon/5)$. Then $0<|x-2|<\delta$ gives
$|x^2-4|<5\delta\leq\epsilon$. The auxiliary restriction controls
the factor that was initially variable.

\paragraph{2. Distinguish displacement and distance.}
A particle has velocity $v(t)=2t-2$ for $0\leq t\leq2$. Find both totals.

\paragraph{Solution.}
Displacement is $[t^2-2t]_0^2=0$. Velocity changes sign at $t=1$.
Distance is $\int_0^1(2-2t)\,dt+\int_1^2(2t-2)\,dt=1+1=2$.
The particle returns to its starting position after traveling two length units.

\paragraph{3. Give an explicit truncation guarantee.}
Approximate $1/(1-x)$ at $x=1/2$ by $\sum_{k=0}^n x^k$. How large
must $n$ be to make the error strictly below $10^{-3}$?

\paragraph{Solution.}
The finite geometric identity gives remainder $x^{n+1}/(1-x)$, hence error
$2^{-n}$ at $x=1/2$. Since $2^{-9}=1/512>10^{-3}$ and
$2^{-10}=1/1024<10^{-3}$, $n=10$ suffices and is minimal.
The approximation uses eleven terms. Continue to
Chapter~\ref{learn:multivariable} to replace a scalar derivative by a linear map.

\chapter{Several variables: local maps and conservation}
\label{learn:multivariable}

\begin{learningcard}{Prerequisites and question}
Use derivatives and integrals from Chapter~\ref{learn:calculus}, and matrix
maps and dot products from Chapter~\ref{learn:linear_algebra}. A temperature
field depends on position, while a moving fluid carries a direction at every
position. How can local rates describe a changing quantity spread across a
region? We work first in Euclidean space with smooth functions and regular
boundaries; singular fields require additional analysis.
\end{learningcard}

\section{A partial derivative controls one input}
A scalar field assigns one real value to each point in a region of $\RR^n$.
For $f(x,y)$, the partial derivative with respect to $x$ is
\[
 f_x(a,b)=\lim_{h\to0}\frac{f(a+h,b)-f(a,b)}h,
\]
when the limit exists. The second input stays fixed. The derivative $f_y$
is defined by varying only the second input. For
$f(x,y)=x^2+3xy+y^2$, ordinary differentiation gives
$f_x=2x+3y$ and $f_y=3x+2y$.

The gradient is the vector $\nabla f=(f_x,f_y)$ in Cartesian coordinates.
If $f$ is differentiable at a point $p$, then
\[
 f(p+h)=f(p)+\nabla f(p)\cdot h+r(h),\qquad
 \frac{|r(h)|}{\|h\|}\longrightarrow0\quad(h\to0).
\]
This condition controls all directions together. The existence of partial
derivatives at one point alone provides weaker information. For example,
set $g(x,y)=xy/(x^2+y^2)$ away from the origin and $g(0,0)=0$.
Both coordinate-axis restrictions vanish, so both partials at the origin
are zero. Along $y=x\ne0$, however, $g=1/2$, so $g$ is discontinuous
at the origin and therefore fails to be differentiable. Continuous first
partials in a neighborhood are a sufficient condition for differentiability.

\section{Worked example: choose a direction of change}
For the polynomial above, $\nabla f(1,0)=(2,3)$. Along a unit direction
$u$, the directional derivative is $\nabla f(1,0)\cdot u$.
The Cauchy--Schwarz inequality bounds its magnitude by $\sqrt{13}$.
For a unit vector $u$, the needed bound follows directly from
$0\leq\|v-(v\cdot u)u\|^2=\|v\|^2-(v\cdot u)^2$;
therefore $|v\cdot u|\leq\|v\|$.
The largest increase occurs at $u=(2,3)/\sqrt{13}$ and equals
$\sqrt{13}$. The direction $(3,-2)/\sqrt{13}$ has zero first-order
change because its dot product with the gradient vanishes.

For an actual displacement $h=(0.01,-0.02)$, the linear prediction is
$2(0.01)+3(-0.02)=-0.04$. Expanding the polynomial shows an exact
additional term $h_x^2+3h_xh_y+h_y^2=-0.0001$, so
$f(1.01,-0.02)-f(1,0)=-0.0401$. The calculation measures the error
instead of treating a local approximation as an identity.

\section{The Jacobian organizes several outputs}
For $F\colon\RR^n\to\RR^m$, its Jacobian matrix has entries
$(DF)_{ij}=\partial F_i/\partial x_j$. Differentiability means
$F(p+h)=F(p)+DF(p)h+r(h)$ with $\|r(h)\|/\|h\|\to0$.
The chain rule becomes matrix multiplication:
\[
 D(G\circ F)(p)=DG(F(p))\,DF(p).
\]
The dimensions enforce the order: an $n$-component displacement enters
$DF$ first, producing an $m$-component displacement for $DG$.

Polar coordinates provide a concrete example:
$F(r,\theta)=(r\cos\theta,r\sin\theta)$. Then
\[
 DF=\begin{pmatrix}\cos\theta&-r\sin\theta\\
 \sin\theta&r\cos\theta\end{pmatrix},\qquad \det DF=r.
\]
The determinant measures signed local area scaling; its absolute value
enters an area integral. Away from $r=0$, a small coordinate rectangle
with sides $dr,d\theta$ therefore has area approximately $r\,dr\,d\theta$.
Coordinate degeneracy at the origin explains why a single global inverse
polar chart requires care.

\begin{learningcard}{Historical situation: notation as shared equipment}
Edwin Bidwell Wilson's 1901 \emph{Vector Analysis}, based on J. Willard
Gibbs's lectures, explicitly presents a textbook for students of mathematics
and physics. Its chapters on differentiation and integration show vector
methods being organized for instruction. The title page documents the
collaboration and pedagogical setting; the chapter organization supplies
evidence of a teachable framework rather than a claim of sole invention.
\footnote{E. B. Wilson, \emph{Vector Analysis: A Text-Book for the Use of
Students of Mathematics and Physics, Founded upon the Lectures of
J. Willard Gibbs} (1901), title page and chapters on differentiation and
integration; digitized text:
\url{https://archive.org/details/vectoranalysiste00gibbiala}.}
Gibbs's preface describes a privately circulated pamphlet printed in
1881--1884 and requests for a fuller public treatment. He then identifies
Wilson as a listener to his 1899--1900 course who undertook the textbook.
Wilson's preface explains his decision to distribute applications through
the exposition as the necessary mathematics became available. These
documented choices show how lectures, limited circulation, and student
needs shaped the published teaching form. Our Jacobian notation and
examples give a modern route through that subject.
\end{learningcard}

\section{Worked example: integration needs the area factor}
The double integral $\iint_D f\,dA$ accumulates values times small areas
and is defined by a limiting sum under suitable integrability assumptions.
For a continuous function on a rectangle, iterated one-variable integration
computes the double integral. Changing variables by a continuously
differentiable one-to-one map with a nonzero Jacobian determinant on the
interior introduces $|\det DF|$.

Compute the integral of $x^2+y^2$ over the unit disk. Polar coordinates
cover the interior apart from the origin and a seam of area zero, which
leave the integral unchanged. The result is
\[
 \iint_{x^2+y^2\leq1}(x^2+y^2)\,dA
 =\int_0^{2\pi}\int_0^1r^2r\,dr\,d\theta
 =2\pi\left[\frac{r^4}{4}\right]_0^1=\frac\pi2.
\]
Dividing by the disk's area $\pi$ gives the mean squared distance from
the center, $1/2$. Omitting the Jacobian factor would count equal radial
widths as equal areas despite their different circumferences.

\section{Divergence turns flux into a local balance}
A vector field $J=(J_1,J_2,J_3)$ assigns a vector to each point. Its
divergence is $\nabla\cdot J=\partial_xJ_1+\partial_yJ_2+\partial_zJ_3$.
For a continuously differentiable field on a neighborhood of a bounded region
$V$ with piecewise smooth boundary, the divergence theorem states
\[
 \int_V\nabla\cdot J\,dV=\int_{\partial V}J\cdot n\,dS,
\]
where $n$ is the outward unit normal and $dS$ is surface area. The right
side is outward flux. On a small rectangular box, opposite-face differences
in each component produce the corresponding partial derivative times volume;
summing boxes cancels shared interior faces. This cancellation explains the
theorem's structure; taking the limit rigorously requires the stated regularity.

Let $\rho(x,t)$ be density, $J(x,t)$ its outward transport flux, and
$s(x,t)$ a local production rate per volume. For a fixed region,
\[
 \frac{d}{dt}\int_V\rho\,dV=-\int_{\partial V}J\cdot n\,dS+\int_Vs\,dV.
\]
Assuming sufficient smoothness to move the derivative inside the integral
and applying the divergence theorem gives
$\int_V(\partial_t\rho+\nabla\cdot J-s)\,dV=0$.
If the continuous integrand integrates to zero on every sufficiently small
region, it vanishes pointwise; otherwise a neighborhood of a strictly
positive or negative value would have a nonzero integral. Hence
$\partial_t\rho+\nabla\cdot J=s$. With $s=0$ and zero boundary flux,
total mass is constant. Chapter~\ref{learn:dynamics} uses this balance to
derive a heat equation and a numerical conservation check.

\section{Exercises with full solutions}
\paragraph{1. Differentiate along a trajectory.}
For $f(x,y)=x^2+y^2$ and $(x(t),y(t))=(\cos t,\sin t)$, find $df/dt$.

\paragraph{Solution.}
The chain rule gives $df/dt=2x x'+2y y'=2\cos t(-\sin t)+
2\sin t\cos t=0$. Direct substitution gives $f=1$, confirming the
result independently. The gradient is radial and the trajectory's velocity
is tangent to the circle, so their dot product vanishes.

\paragraph{2. Verify a divergence theorem example.}
For $J=(x,y,z)$ on the unit cube $[0,1]^3$, compute both sides.

\paragraph{Solution.}
The divergence is three, so its volume integral is three. On each face
$x=0$, $y=0$, or $z=0$, the normal component is zero. On each opposite
face the outward normal component is one and the area is one. The total
outward flux is therefore $1+1+1=3$, matching the volume integral.

\paragraph{3. Identify the conservation boundary.}
In a one-dimensional interval $[0,L]$, suppose $\rho_t+J_x=0$.
Express the mass derivative and identify a sufficient boundary condition.

\paragraph{Solution.}
Integrating and using the fundamental theorem yields
$M'(t)=-[J]_0^L=J(0,t)-J(L,t)$, where $M=\int_0^L\rho\,dx$.
Equal endpoint fluxes guarantee $M'=0$; zero flux at both endpoints is one
such condition. Periodic flux is another. The local equation alone allows
mass to leave through the boundary, so the boundary convention belongs in
any global conservation claim.

\chapter{Dynamics: equations, modes, and trustworthy computation}
\label{learn:dynamics}

\begin{learningcard}{Prerequisites and guiding experiment}
Use derivatives and integrals from Chapter~\ref{learn:calculus}, eigenvectors
from Chapter~\ref{learn:linear_algebra}, and conservation from
Chapter~\ref{learn:multivariable}. Imagine a warm bar cooling toward a
uniform temperature. We need an evolution equation, initial and boundary
data, a way to solve a controlled example, and a check that a computer's
approximation respects the mechanism.
\end{learningcard}

\section{An equation of motion needs initial data}
An ordinary differential equation (ODE) relates an unknown function of one
variable to its derivatives. For $y'=f(t,y)$, an initial condition
$y(t_0)=y_0$ chooses a trajectory when an existence and uniqueness theorem
applies. Continuity in $t$ and local Lipschitz continuity in $y$ near the
initial point provide a standard local guarantee. Lipschitz continuity means
$|f(t,u)-f(t,v)|\leq L|u-v|$ locally for a fixed finite constant $L$.
The guarantee is local: solutions can grow without bound in finite time,
as $y'=y^2$, $y(0)=1$, with $y=1/(1-t)$ illustrates.

\section{Worked example: exponential relaxation}
Suppose excess temperature $y$ obeys $y'=-ky$, where $k>0$ has units of
inverse time. Multiplying by $e^{kt}$ gives
$(e^{kt}y)'=e^{kt}(y'+ky)=0$. Therefore
\[
 y(t)=y_0e^{-kt}.
\]
This integrating-factor argument includes $y_0=0$, a case division by $y$
would have excluded temporarily. The time to reduce a positive excess by
half is $t_{1/2}=\log2/k$. With $k=0.2$ per minute and $y_0=10$ degrees,
the excess after five minutes is $10/e$ degrees, about $3.679$.
The exponential law is a model assumption. Its success in one regime does
not establish the same coefficient for every material or temperature range.

An equilibrium of an autonomous equation $y'=f(y)$ is a value $y_*$ with
$f(y_*)=0$. It is stable when sufficiently small initial perturbations stay
small for all forward time for which the comparison is made; it is
asymptotically stable when nearby trajectories also approach it. In the
relaxation model, the exact factor $e^{-kt}$ proves both properties at zero.

\section{Worked example: linear stability has directional content}
For a constant matrix system $z'=Az$, an eigenvector $v$ with eigenvalue
$\lambda$ produces the solution $z(t)=ce^{\lambda t}v$. If a basis of
eigenvectors exists, expand the initial vector in that basis and evolve
each coefficient. Consider
\[
 A=\begin{pmatrix}-2&1\\1&-2\end{pmatrix}.
\]
The vectors $(1,1)^T$ and $(1,-1)^T$ have eigenvalues $-1$ and $-3$.
For $z(0)=(2,0)^T$, the decomposition gives
\[
 z(t)=e^{-t}(1,1)^T+e^{-3t}(1,-1)^T.
\]
Both modes decay, and the slower symmetric mode dominates at late times.
A matrix with an eigenvalue of positive real part instead has a growing
mode. In finite dimensions, every eigenvalue having negative real part
ensures asymptotic stability even when generalized eigenvectors are needed.
Eigenvalues on the imaginary axis require more information.

For a smooth nonlinear system, write $z=z_*+\eta$ near an equilibrium.
Differentiability gives $\eta'=Df(z_*)\eta+o(\|\eta\|)$.
Strictly negative real parts imply local asymptotic stability under the
standard linearization theorem. A zero eigenvalue makes that test
inconclusive: $y'=-y^3$ attracts nearby trajectories, while $y'=y^3$
repels them, and both derivatives at zero vanish.

\begin{learningcard}{Historical situation: heat and a research program}
Fourier's \emph{Th\'eorie analytique de la chaleur} (1822) opens with a
preliminary discourse placing heat among mathematical investigations of
natural phenomena. The text develops equations and trigonometric expansions
for thermal problems. Its introductory account also describes the submission
of work to the Institut. That setting makes scientific institutions part of
the documented story of circulation and assessment.\footnote{Joseph Fourier,
\emph{The Analytical Theory of Heat}, translated by Alexander Freeman
(1878), Preliminary Discourse and Chapter I; University of Michigan
Historical Mathematics Collection:
\url{https://quod.lib.umich.edu/u/umhistmath/abr1558.0001.001}.}
The source documents a program and its presentation. Modern convergence and
solution theorems specify which Fourier manipulations are justified.
\end{learningcard}

\section{From a conservation law to the heat equation}
A partial differential equation (PDE) involves an unknown function of several
variables and its partial derivatives. In one space dimension let $u(x,t)$
be temperature, $c>0$ constant volumetric heat capacity, and $J$ heat flux.
Local energy conservation without sources says $cu_t+J_x=0$. Fourier's
constitutive law $J=-K u_x$ assumes a constant conductivity $K>0$ and flux
opposite the temperature gradient. Combining the two gives
\[
 u_t=\kappa u_{xx},\qquad\kappa=K/c>0.
\]
Here $\kappa$ has units of length squared per time. Conservation supplies
a balance; the constitutive law supplies material behavior. They are separate
assumptions that experiments can examine.

On $0\leq x\leq L$, impose zero endpoint temperature and initial profile
$u(x,0)=\sin(\pi x/L)$. Direct differentiation verifies the exact solution
\[
 u(x,t)=e^{-\kappa(\pi/L)^2t}\sin(\pi x/L).
\]
More generally, a finite linear combination of sine modes evolves by
multiplying mode $n$ by $e^{-\kappa(n\pi/L)^2t}$. Higher spatial frequencies
decay faster. Infinite expansions require convergence and regularity
conditions; an arbitrary interchange of a sum and a derivative needs proof.
For sufficiently smooth solutions with zero endpoint temperatures,
integration by parts gives
\[
 \frac{d}{dt}\frac12\int_0^L u^2\,dx
 =\kappa\int_0^Luu_{xx}\,dx=-\kappa\int_0^L u_x^2\,dx\leq0.
\]
The boundary term vanishes because $u$ vanishes there. This squared-norm
quantity supplies a stability check; total heat may leave through the ends.

\section{Stable dynamics and unstable time steps}
Forward Euler replaces $y'$ by a forward difference. For $y'=-ky$,
\[
 y_{m+1}=(1-k\Delta t)y_m.
\]
The exact solution decays for every positive time step, but the numerical
factor has magnitude at most one only when $0\leq k\Delta t\leq2$.
Strict decay requires $0<k\Delta t<2$; the endpoint two gives undamped
sign alternation. Nonnegativity for positive initial data requires the stronger
condition $k\Delta t\leq1$. At equality the iterate is zero; strict positivity requires $k\Delta t<1$.
Accuracy, stability, and sign preservation are distinct.

With equally spaced spatial points and time levels, the centered-space,
forward-time heat scheme is
\[
 u_j^{m+1}=u_j^m+r(u_{j+1}^m-2u_j^m+u_{j-1}^m),\qquad
 r=\frac{\kappa\Delta t}{\Delta x^2}.
\]
For periodic data, insert a mode $e^{ij\theta}$, where $i^2=-1$.
Its amplification factor is
$G(\theta)=1-4r\sin^2(\theta/2)$. The condition $0\leq r\leq1/2$
guarantees $|G|\leq1$ for every frequency; frequencies near $\pi$ expose
failure beyond that bound. Spatial refinement by a factor two therefore
requires a time step at most one quarter as large under this sufficient
stability rule.

A computation should record initial data, boundary rules, coefficients,
mesh widths, and time steps. Compare a known mode with its exact solution,
repeat on refined grids, and measure conservation under a boundary condition
that actually conserves the quantity. These checks test the discretization
and implementation; establishing a physical model additionally needs
observations and uncertainty estimates.

\section{Exercises with full solutions}
\paragraph{1. Test two Euler time steps.}
For $y'=-2y$, compare $\Delta t=0.4$ and $\Delta t=1.1$.

\paragraph{Solution.}
The amplification factors are $1-2(0.4)=0.2$ and
$1-2(1.1)=-1.2$. The first gives $y_m=0.2^my_0$, decaying with
preserved sign. The second gives $y_m=(-1.2)^my_0$, alternating and
growing in magnitude for nonzero initial data. The exact solution
$y_0e^{-2t}$ decays in both comparisons, identifying the growth as numerical.

\paragraph{2. Follow two heat modes.}
On $[0,\pi]$ with zero endpoints and $\kappa=1$, evolve
$u(x,0)=\sin x+2\sin3x$.

\paragraph{Solution.}
Each mode solves the equation separately, and linearity gives
$u(x,t)=e^{-t}\sin x+2e^{-9t}\sin3x$. Substituting $t=0$ verifies
the initial condition; both sines vanish at the endpoints. Two spatial
derivatives multiply the modes by $-1$ and $-9$, matching the time
derivatives. The coefficient ratio is $2e^{-8t}$, so the higher-frequency
contribution decays faster.

\paragraph{3. Prove discrete mass conservation.}
For the periodic heat scheme with $N$ points, prove preservation of
$M^m=\Delta x\sum_{j=0}^{N-1}u_j^m$ in exact arithmetic.

\paragraph{Solution.}
Summing the update gives
$M^{m+1}-M^m=r\Delta x\sum_j(u_{j+1}^m-2u_j^m+u_{j-1}^m)$.
Periodic indexing makes each shifted sum equal to $\sum_j u_j^m$,
so the difference is zero. Floating-point addition can introduce rounding
drift, which a program should measure. Conservation alone also leaves the
stability question open: the sum cancellation holds even for a time step
that amplifies oscillating modes.

\chapter{Probability: uncertainty, evidence, and comparison}
\label{learn:probability}

\begin{learningcard}{Prerequisites and question}
Bring fractions, sums, functions, and the logic of Chapter~\ref{learn:proof}.
Calculus becomes useful for continuous distributions, but our main examples
have finite sample spaces. A positive diagnostic result sounds decisive.
How should its meaning depend on the condition's prevalence and the test's
errors? The same reasoning helps interpret a simulation's apparent success.
\end{learningcard}

\section{Specify the experiment before counting}
A sample space $\Omega$ is the set of possible outcomes. An event is a
subset of outcomes whose occurrence we can ask about. On a finite space,
assign probabilities $p(\omega)\geq0$ with
$\sum_{\omega\in\Omega}p(\omega)=1$; the probability of an event is the
sum over its outcomes. Equal probabilities require a modeling assumption,
such as an ideal fair die. Counting favorable outcomes divided by all outcomes
works only under that assumption.

A random variable $X$ is a real-valued function of the outcome. Its expectation
and variance are
\[
 \mathbb E[X]=\sum_{\omega}p(\omega)X(\omega),\qquad
 \operatorname{Var}(X)=\mathbb E[(X-\mathbb E[X])^2].
\]
Expectation is a probability-weighted mean, and variance measures squared
spread about that mean. Expanding the square gives
$\operatorname{Var}(X)=\mathbb E[X^2]-(\mathbb E[X])^2$.
Standard deviation is the nonnegative square root of variance and has the
same units as $X$.

\section{Worked example: a Bernoulli observation}
Let $X$ equal one when a trial succeeds and zero otherwise, with success
probability $p$. Then $\mathbb E[X]=p$, and $X^2=X$ gives
$\operatorname{Var}(X)=p-p^2=p(1-p)$. For independent identically
distributed trials $X_1,\ldots,X_n$, their sample mean
$\overline X=n^{-1}\sum_jX_j$ has expectation $p$. Independence makes
cross covariances zero, so
\[
 \operatorname{Var}(\overline X)=\frac{p(1-p)}n.
\]
The standard deviation of the mean therefore scales as $1/\sqrt n$.
Four times as many independent observations halve it. Repeating the same
stored observation four times creates dependence and supplies no such gain.

More generally, covariance is
$\operatorname{Cov}(X,Y)=\mathbb E[(X-\mathbb E[X])(Y-\mathbb E[Y])]$.
The variance of a sum includes twice the pairwise covariances. Correlation
normalizes covariance by the two positive standard deviations. Independence
implies zero covariance when the moments exist; the converse can fail.
For a uniformly distributed $X$ on $\{-1,0,1\}$, set $Y=X^2$.
Then $\mathbb E[X]=\mathbb E[X^3]=0$, so covariance vanishes, although
$Y$ is determined by $X$.

\section{Worked example: condition on the result you observed}
For events $A,B$ with $P(B)>0$, conditional probability is
$P(A\mid B)=P(A\cap B)/P(B)$. Applying this definition in both orders
yields Bayes's rule:
\[
 P(A\mid B)=\frac{P(B\mid A)P(A)}{P(B)}.
\]
Let $D$ mean a condition is present, and $+$ a positive test. Assume a
population prevalence $P(D)=0.01$, sensitivity $P(+\mid D)=0.99$,
and false-positive probability $P(+\mid D^c)=0.05$. These are hypothetical
teaching parameters, rather than clinical advice about a real test.
Partitioning by condition status gives
\[
 P(+)=0.99(0.01)+0.05(0.99)=0.0594,
 \qquad P(D\mid+)=\frac{0.0099}{0.0594}=\frac16.
\]
A hypothetical population of 10,000 makes the denominator visible: among
100 people with the condition we expect 99 positive results; among 9,900
others we expect 495. Only 99 of the expected 594 positives come from the
condition group. Expected counts may be nonintegers in other examples;
they describe averages under the model rather than guaranteed outcomes.

\begin{figure}[htbp]
  \centering
  \begin{tikzpicture}[x=1cm,y=1cm,>=Stealth,font=\small,line width=0.8pt]
  \node[align=center] (population) at (0,0) {Population};
  \node (condition) at (3.5,1.8) {$D$};
  \node (other) at (3.5,-1.8) {$D^c$};
  \node[anchor=west] (truepositive) at (7,2.7) {$+$: $P(D\cap+)=99/10000$};
  \node[anchor=west] (falsenegative) at (7,0.9) {$-$: $P(D\cap-)=1/10000$};
  \node[anchor=west] (falsepositive) at (7,-0.9) {$+$: $P(D^c\cap+)=495/10000$};
  \node[anchor=west] (truenegative) at (7,-2.7) {$-$: $P(D^c\cap-)=9405/10000$};
  \draw[->] (population)--(condition) node[midway,above,sloped] {$P(D)=1/100$};
  \draw[->] (population)--(other) node[midway,below,sloped] {$P(D^c)=99/100$};
  \draw[->,line width=1.2pt] (condition)--(truepositive.west)
    node[midway,above,sloped] {$P(+\mid D)=99/100$};
  \draw[->,dashed] (condition)--(falsenegative.west)
    node[midway,below,sloped] {$P(-\mid D)=1/100$};
  \draw[->,line width=1.2pt] (other)--(falsepositive.west)
    node[midway,above,sloped] {$P(+\mid D^c)=5/100$};
  \draw[->,dashed] (other)--(truenegative.west)
    node[midway,below,sloped] {$P(-\mid D^c)=95/100$};
  \node[align=center] at (5.8,-3.7)
    {Keep the two positive leaves, then normalize:\\[3pt]
     $\displaystyle P(D\mid+)=\frac{99/10000}{(99+495)/10000}=\frac16$.};
\end{tikzpicture}

  \caption{Read each path from left to right and multiply its conditional branch probabilities to obtain the joint probability at the leaf. Restrict attention to the two solid positive-result branches, add their leaf probabilities, and normalize the condition-positive leaf by that total.}
  \label{fig:learn:bayes_tree}
\end{figure}

\begin{learningcard}{Historical situation: a paper communicated by another reader}
Thomas Bayes's essay appeared in 1763 after his death, communicated by
Richard Price to the Royal Society. The publication's introductory letter
makes Price's role in presenting the work part of the surviving evidence.
The essay investigates inverse-probability reasoning; our diagnostic example
is a modern application rather than a historical episode from the essay.
\footnote{Thomas Bayes, ``An Essay towards solving a Problem in the
Doctrine of Chances,'' communicated by Richard Price,
\emph{Philosophical Transactions} 53 (1763), pp.~370--418, introductory
letter and propositions:
\url{https://doi.org/10.1098/rstl.1763.0053}.}
A named theorem can conceal the editorial and institutional work through
which a mathematical argument reaches its readers.
\end{learningcard}

\section{An interval is a procedure with assumptions}
Suppose observations are independent normal random variables with unknown
mean $\mu$ and known positive standard deviation $\sigma$. Their mean is
normal with standard deviation $\sigma/\sqrt n$. If $z_{0.975}$ denotes
the 97.5th percentile of the standard normal distribution, approximately
$1.96$, then the interval
\[
 \left[\overline X-z_{0.975}\frac{\sigma}{\sqrt n},\,
       \overline X+z_{0.975}\frac{\sigma}{\sqrt n}\right]
\]
has 95 percent coverage: over repeated samples under these assumptions,
95 percent of intervals contain the fixed $\mu$. A particular observed
interval either contains that fixed value or misses it. A Bayesian credible
interval instead assigns posterior probability to parameter values under a
specified prior and likelihood.

Unknown variance, nonnormal small samples, dependence, selection, and repeated
comparisons change the calculation or its interpretation. For independent
normal observations with unknown variance, a Student $t$ interval uses the
sample standard deviation and $n-1$ degrees of freedom. Large-sample normal
approximations require justification. Running many tests and reporting only
the most striking interval changes the procedure's coverage claims.

\section{A control separates explanations}
Association is a property of a joint distribution. A causal claim asks what
would change under an intervention. Both temperature and ice-cream sales
can correlate with swimming incidents; a shared seasonal cause can generate
association without the proposed direct mechanism. A randomized comparison
balances causes in expectation under its design assumptions. Observational
causal inference needs additional assumptions about confounding and selection.

For simulations, matched initial conditions and recorded parameter changes
help isolate the mechanism being compared. Independent seed clusters matter
when outputs within one run share a trajectory. Report the difference between
models with uncertainty, rather than interpreting one model's attractive
picture as evidence of a specific cause. Chapter~\ref{learn:numerical}
adds implementation error to these statistical questions.

\section{Exercises with full solutions}
\paragraph{1. Change the population.}
Keep the diagnostic sensitivity and false-positive probability above, but
set prevalence to $0.10$. Find $P(D\mid+)$.

\paragraph{Solution.}
The positive probability is $0.99(0.10)+0.05(0.90)=0.144$.
The joint probability of the condition and a positive result is $0.099$,
so $P(D\mid+)=0.099/0.144=11/16=0.6875$.
The same test characteristics give a different posterior probability because
the population entering the calculation changed.

\paragraph{2. Compute an uncertainty scale.}
Independent Bernoulli trials have $p=1/2$. Find the standard deviation of
the sample mean for $n=100$ and $n=400$.

\paragraph{Solution.}
The formula gives $\sqrt{(1/2)(1/2)/100}=0.05$ and
$\sqrt{(1/2)(1/2)/400}=0.025$. Both numbers describe the sampling
spread under known $p=1/2$; estimating an unknown $p$ adds an inferential
step. Fourfold sample size halves the standard deviation.

\paragraph{3. Check whether two events are independent.}
For a fair die, let $A$ be an even result and $B$ a result at least four.

\paragraph{Solution.}
$P(A)=3/6=1/2$ and $P(B)=3/6=1/2$. Their intersection is
$\{4,6\}$, so $P(A\cap B)=1/3$, which differs from $P(A)P(B)=1/4$.
The events are dependent. Equivalently,
$P(A\mid B)=2/3\ne1/2=P(A)$.

\chapter{Numerical reasoning: approximation with an error budget}
\label{learn:numerical}

\begin{learningcard}{Prerequisites and task}
Use Taylor bounds from Chapter~\ref{learn:calculus}, matrix maps from
Chapter~\ref{learn:linear_algebra}, and stability from
Chapter~\ref{learn:dynamics}. A numerical answer is useful when its precision
matches the question and its errors are understood. We ask how to distinguish
an accurate-looking output from an approximation that survives refinement
and independent checks.
\end{learningcard}

\section{Three errors can coexist}
Modeling error measures the difference between a chosen mathematical model
and the target phenomenon. Discretization or truncation error comes from
replacing an infinite or continuous operation by a finite one. Roundoff error
comes from representing and manipulating numbers with finite precision.
Changing the mesh can reduce discretization error while leaving a wrong model
unchanged. Exact rational arithmetic can remove roundoff for some calculations
while leaving truncation error intact.

A problem is well-conditioned when small perturbations of its input produce
small relative changes in the requested output, with scales specified.
For a differentiable scalar function away from zero input and output, the
relative condition number is $|xf'(x)/f(x)|$. An algorithm's stability
instead concerns how its errors behave while solving the problem. A stable
algorithm can still produce an uncertain answer for an ill-conditioned input.

\section{Worked example: a smaller difference can lose information}
Approximate $f'(x)$ by
$D_hf(x)=[f(x+h)-f(x)]/h$ for $h\ne0$. Taylor's theorem gives
$D_hf(x)=f'(x)+hf''(\xi)/2$ for a suitable intermediate point when the
required smoothness holds. If $|f''|\leq M$, truncation error is at most
$M|h|/2$.

Suppose each function evaluation has absolute error at most $\eta$.
Subtracting the evaluations and dividing by $h$ adds an error at most
$2\eta/|h|$, before accounting for the arithmetic operations themselves.
The combined illustrative bound is
\[
 |D_h^{\mathrm{computed}}f(x)-f'(x)|
 \leq \frac{M|h|}{2}+\frac{2\eta}{|h|}.
\]
For $M,\eta>0$, balancing these two terms suggests
$|h|=2\sqrt{\eta/M}$. Arbitrarily shrinking $h$ amplifies evaluation
error. This formula depends on the absolute-error model; a real program
must account for scales and the actual evaluation routine.

For $f(x)=x^2$ at $x=1$, exact algebra gives $D_hf(1)=2+h$.
At $h=0.01$, the error is exactly $0.01$. A floating-point implementation
may eventually round $1+h$ to one, producing a zero numerator despite the
true derivative two. An exact symbolic identity provides a better oracle
than merely comparing two nearby floating-point outputs.

\section{Worked example: integration with a computable remainder}
The composite trapezoidal rule on $[a,b]$ with $n$ equal pieces of width
$h=(b-a)/n$ is
\[
 T_n=h\left(\frac{f(a)+f(b)}2+\sum_{j=1}^{n-1}f(a+jh)\right).
\]
For $f$ with continuous second derivative bounded by $M$, its error obeys
$|T_n-\int_a^bf(x)\,dx|\leq(b-a)Mh^2/12$.
One obtains the order by integrating the local linear-interpolation error
on each piece: a local error of order $h^3$ accumulated across $n$ pieces
gives a global error of order $h^2$.

For $f(x)=x^2$ on $[0,1]$, use
$\sum_{j=1}^{n-1}j^2=(n-1)n(2n-1)/6$ to derive
\[
 T_n=\frac1{2n}+\frac{(n-1)n(2n-1)}{6n^3}
 =\frac13+\frac1{6n^2}.
\]
The exact integral is $1/3$, so the error is positive and equals
$1/(6n^2)$. Since $M=2$, the general bound is attained. Doubling $n$
reduces the error by four. The extrapolated quantity
$(4T_{2n}-T_n)/3$ equals $1/3$ exactly for this quadratic because the
leading error formula is exact. For other smooth functions, extrapolation
requires the appropriate error expansion.

\begin{learningcard}{Historical situation: reliable matrix computation}
Turing's 1948 paper ``Rounding-off Errors in Matrix Processes'' studies
error in numerical matrix operations. The paper's subject is a concrete
reminder that computation requires analysis of the procedure as well as
algebraic correctness of the equations. Its publication documents an early
analysis of such errors; it supplies neither the origin story of all
numerical analysis nor a guarantee for a modern implementation.
\footnote{Alan M. Turing, ``Rounding-off Errors in Matrix Processes,''
\emph{Quarterly Journal of Mechanics and Applied Mathematics} 1 (1948),
pp.~287--308: \url{https://doi.org/10.1093/qjmam/1.1.287}.}
The finite-difference experiment above uses a modern explicit error budget.
\end{learningcard}

\section{Iteration: a residual and an error answer different questions}
To solve $Ax=b$, consider $x_{m+1}=Bx_m+c$. If a fixed point $x_*$
exists, its error obeys $e_{m+1}=Be_m$. Any induced norm with $\|B\|<1$
gives $\|e_m\|\leq\|B\|^m\|e_0\|$, proving convergence. In finite
dimensions, spectral radius below one is the general constant-matrix criterion
for convergence from every initial error, although the chosen norm can show
transient growth before decay.

For the system $2x+y=3$, $x+2y=3$, simultaneous updates give
$x_{m+1}=(3-y_m)/2$, $y_{m+1}=(3-x_m)/2$. The error matrix is
\[
 B=\begin{pmatrix}0&-1/2\\-1/2&0\end{pmatrix}.
\]
Its Euclidean operator norm is $1/2$. Starting at $(0,0)$ produces
$(1.5,1.5)$, $(0.75,0.75)$, $(1.125,1.125)$, approaching $(1,1)$.
The residual is $r=b-A\widehat x$. If $A$ is invertible, the actual
error satisfies $\widehat x-x_*=-A^{-1}r$, so
$\|\widehat x-x_*\|\leq\|A^{-1}\|\|r\|$. A small residual establishes
a small solution error only when the inverse scale is controlled.

\section{A convergence study needs an oracle and a regime}
For a time-dependent scheme, consistency asks whether the discrete equation
approximates the differential equation as the spacings shrink. Stability
controls growth of perturbations. Convergence asks whether the computed
solution approaches the desired exact solution. The familiar equivalence
between consistency plus stability and convergence has specified hypotheses
for linear well-posed problems; it is not a universal statement about every
nonlinear simulation.

Conservation gives another independent test. As
Chapter~\ref{learn:dynamics} shows, the periodic explicit heat scheme
conserves its discrete sum even at an unstable time step. A mass check can
pass while oscillating errors grow. Conversely, a nonconserving boundary
condition can legitimately change total mass.

Record a normed error against a known solution on successive grids, estimate
the ratio, and verify that boundary and time-step conventions refine together.
A ratio near the expected value supports the observed regime. Roundoff,
singularities, or unresolved scales can prevent that regime from appearing.
Chapter~\ref{learn:probability} supplies uncertainty analysis when random
samples enter the computation; that uncertainty differs from deterministic
mesh error.

\section{Exercises with full solutions}
\paragraph{1. Choose an integration resolution.}
For the quadratic trapezoidal example, find the smallest positive integer
$n$ guaranteeing error at most $10^{-4}$.

\paragraph{Solution.}
The exact error is $1/(6n^2)$. The condition requires
$n^2\geq10000/6$, hence $n\geq41$. At $n=40$ the error is
$1/9600>10^{-4}$, whereas at $n=41$ it is $1/10086<10^{-4}$.
Thus 41 pieces suffice and are minimal for this specified rule and integrand.

\paragraph{2. Bound an iteration count.}
In the two-variable example, how many updates guarantee Euclidean error
below $0.01$ from $(0,0)$?

\paragraph{Solution.}
The initial error norm is $\sqrt2$ and the bound is $2^{-m}\sqrt2$.
At $m=7$ this is approximately $0.01105$; at $m=8$ it is approximately
$0.00552$. Eight updates suffice. The initial error lies in an eigenvector
direction of magnitude-one-half eigenvalue, so the bound is exact here.

\paragraph{3. Interpret a refinement table.}
Errors at spacings $h,h/2,h/4$ are $0.04,0.01,0.0025$. Estimate the
observed order and state the conclusion's boundary.

\paragraph{Solution.}
For $E(h)\approx Ch^p$, the ratio $E(h)/E(h/2)\approx2^p$.
Both measured ratios are four, giving $p=\log_2 4=2$.
The three measurements support second-order behavior on those spacings for
the specified test. They alone establish neither every-grid convergence
nor correctness of the physical model.

\chapter{Algebraic structure: symmetry, brackets, and roots}
\label{learn:algebra}

\begin{learningcard}{Prerequisites and destination}
Use proof, vector spaces, matrices, and eigenvectors from
Chapters~\ref{learn:proof} and~\ref{learn:linear_algebra}. The aim is to
recognize what an algebraic classification classifies before attaching a
physical interpretation to it. We move from rotations of a triangle to a
small root system, then identify concrete features of the exceptional
$E_8$ endpoint.
\end{learningcard}

\section{A group records reversible composition}
A group is a set $G$ with an associative operation, an identity element,
and an inverse for every element; the operation stays in $G$. If the
operation also commutes, the group is abelian. The rotations of an
equilateral triangle by multiples of $120$ degrees form a three-element
group. Label them $0,1,2$ and compose by addition modulo three: two
successive $120$-degree turns give label two, and three give the identity.
Rotations by arbitrary angles about one fixed axis similarly compose by
angle addition modulo a full turn.

A permutation is a bijection from a set to itself. For three labeled
vertices, the six permutations form the symmetric group $S_3$ under
composition. Write $(12)$ for swapping labels one and two. With the rightmost
permutation acting first, $(12)(23)$ sends one to two, two to three, and
three to one. Reversing the order gives the inverse cycle. Therefore $S_3$
is nonabelian even though its subgroup of three rotations is abelian.

A group action on a set assigns each group element a transformation so that
identity and composition agree. Specifying an action matters: naming a group
beside a physical system does not define how that group's elements transform
the system's states or observations.

\section{Structures supply different operations}
A field has addition and multiplication satisfying the familiar distributive
and associative laws, commutative multiplication, distinct zero and one, and
a multiplicative inverse for every nonzero element. The real and complex
numbers are fields. A vector space over a field has vector addition and
scalar multiplication; it need not multiply its vectors together. An algebra
over a field is a vector space with an additional bilinear product, meaning
a product linear in each argument separately. Extra adjectives specify
whether that product is associative, commutative, or has an identity.

Real square matrices form an associative algebra under multiplication and a
vector space under addition, but they form a multiplicative group only after
restricting to invertible matrices. All matrices together fail that group
condition because singular matrices have no inverse. Keeping these operations
separate prevents a dimension count from substituting for a structural proof.

\section{Worked example: a bracket measures failure to commute}
For matrices define $[A,B]=AB-BA$. This bracket is bilinear,
$[A,A]=0$, and satisfies the Jacobi identity
\[
 [A,[B,C]]+[B,[C,A]]+[C,[A,B]]=0.
\]
To verify Jacobi, expand each nested bracket: the resulting twelve terms
cancel in pairs using associativity. A vector space with a bilinear bracket
satisfying these identities is a Lie algebra.

Consider the trace-zero matrices
\[
 E=\begin{pmatrix}0&1\\0&0\end{pmatrix},\quad
 F=\begin{pmatrix}0&0\\1&0\end{pmatrix},\quad
 H=\begin{pmatrix}1&0\\0&-1\end{pmatrix}.
\]
Multiplication gives $EF=\operatorname{diag}(1,0)$ and
$FE=\operatorname{diag}(0,1)$, hence $[E,F]=H$. Computing the other
products gives $[H,E]=2E$ and $[H,F]=-2F$. These three matrices span
the trace-zero two-by-two matrices, conventionally denoted
$\mathfrak{sl}_2$ over the chosen real or complex field. The eigenvector
language from linear algebra reappears: under the linear map
$X\mapsto[H,X]$, the vectors $E,F$ have eigenvalues two and minus two.

\begin{learningcard}{Historical situation: classification requires definitions}
Cartan's 1894 dissertation investigates the structure of finite continuous
transformation groups. The historical object is a published structural
investigation, not a claim that every later application was contained in that
work. The dissertation belongs to the development from continuous
transformations toward the Lie-algebra classifications used in the advanced
review \citep{cartan1894}. Our matrix brackets provide a small calculation
through which a reader can see what such structural questions ask.
\end{learningcard}

\section{Worked example: the six roots of \texorpdfstring{$A_2$}{A2}}
Let $e_1,e_2,e_3$ be the standard unit vectors in $\RR^3$. The six vectors
$e_i-e_j$ with $i\ne j$ lie in the plane $x_1+x_2+x_3=0$.
They form the root system called $A_2$. Choose simple roots
$\alpha_1=e_1-e_2$ and $\alpha_2=e_2-e_3$; the remaining positive root
is $\alpha_1+\alpha_2=e_1-e_3$, with negatives supplying the other three.
Both simple roots have squared length two and their dot product is minus one.
The angle between them therefore has cosine $-1/2$.

\begin{figure}[htbp]
  \centering
  \begin{tikzpicture}[x=1.65cm,y=1.65cm,>=Stealth,font=\small,line width=0.8pt]
  \coordinate (first) at ({sqrt(2)},0);
  \coordinate (sum) at ({sqrt(2)/2},{sqrt(6)/2});
  \coordinate (second) at ({-sqrt(2)/2},{sqrt(6)/2});
  \coordinate (minusfirst) at ({-sqrt(2)},0);
  \coordinate (minussum) at ({-sqrt(2)/2},{-sqrt(6)/2});
  \coordinate (minussecond) at ({sqrt(2)/2},{-sqrt(6)/2});
  \draw[dashed] (first)--(sum)--(second)--(minusfirst)--(minussum)--(minussecond)--cycle;
  \foreach \point in {first,sum,second,minusfirst,minussum,minussecond} {
    \draw[->,line width=1.1pt] (0,0)--(\point);
  }
  \node[right] at (first) {$\alpha_1$};
  \node[above right] at (sum) {$\alpha_1+\alpha_2$};
  \node[above left] at (second) {$\alpha_2$};
  \node[left] at (minusfirst) {$-\alpha_1$};
  \node[below left] at (minussum) {$-(\alpha_1+\alpha_2)$};
  \node[below right] at (minussecond) {$-\alpha_2$};
  \draw[->] (0.45,0) arc[start angle=0,end angle=120,radius=0.45];
  \node[fill=white,inner sep=1pt] at (0.54,0.3) {$120^\circ$};
  \node[below] at (0,-0.1) {$0$};
  \node[align=center] at (0,-2)
    {All six arrows have length $\sqrt2$.\\
     $\alpha_1\cdot\alpha_2=-1$; $\cos120^\circ=-1/2$.};
\end{tikzpicture}

  \caption{Read the two simple-root arrows first: their angle is 120 degrees. Their vector sum supplies the upper-right root, and negation supplies the opposite arrows. The picture uses orthonormal coordinates in the root plane; the simple roots themselves are not orthogonal.}
  \label{fig:learn:root_hexagon}
\end{figure}

For nonzero roots, reflection perpendicular to $\alpha$ is
$s_\alpha(v)=v-2(v\cdot\alpha)\alpha/(\alpha\cdot\alpha)$.
Reflection in $\alpha_1=e_1-e_2$ swaps the first two coordinates, so it
permutes the six-root set. The Cartan matrix records relative inner products,
$C_{ij}=2(\alpha_i\cdot\alpha_j)/(\alpha_j\cdot\alpha_j)$; here
\[
 C=\begin{pmatrix}2&-1\\-1&2\end{pmatrix},\qquad\det C=3.
\]
The root lattice $Q$ consists of integer combinations of the simple roots.
The weight lattice $P$ consists of vectors pairing integrally with each
coroot $2\alpha_j/(\alpha_j\cdot\alpha_j)$. Solving those pairings
introduces thirds because
\[
 C^{-1}=\frac13\begin{pmatrix}2&1\\1&2\end{pmatrix}.
\]
For this full-rank finite root system, $P/Q$ has three elements. It is a
lattice quotient rather than an automatic selection rule for a fluid.

\section{A concrete endpoint: the \texorpdfstring{$E_8$}{E8} count and lengths}
A standard coordinate construction in $\RR^8$ has two kinds of roots:
\[
 (\pm1,\pm1,0,0,0,0,0,0)\text{ with all position choices},
 \qquad \tfrac12(\pm1,\ldots,\pm1)\text{ with even minus parity}.
\]
There are $\binom82\cdot4=112$ roots of the first kind. Of the $2^8$
sign patterns of the second kind, exactly half have even minus parity,
giving 128. The total is 240. Every first-kind root has squared length
$1+1=2$; every second-kind root has squared length $8/4=2$.
The two families therefore share the same length $\sqrt2$.

The $E_8$ Cartan diagram is a tree with one central vertex and three arms
of lengths one, two, and four. Its matrix has diagonal entries two and
entries minus one along edges. An arm of length $m$ has determinant $m+1$;
this follows from the tridiagonal recurrence $D_m=2D_{m-1}-D_{m-2}$,
$D_0=1,D_1=2$. The inverse entry at its attaching end is $m/(m+1)$,
by the cofactor formula. Eliminating the arms thus gives
\[
 \det C=2\cdot3\cdot5
 \left(2-\frac12-\frac23-\frac45\right)=1.
\]
The corresponding root and weight lattices coincide. Classification identifies
this root system with the rank-eight complex simple Lie algebra $E_8$,
whose dimension is eight Cartan directions plus 240 one-dimensional root
spaces, totaling 248. Proving that classification and constructing all
representations belong to further Lie theory. The finite calculations here
verify the stated count, lengths, and determinant, rather than that entire
classification theorem \citep{kac1990infinite}.

\section{Exercises with full solutions}
\paragraph{1. Compute an inverse in the triangle rotation group.}
Find the inverse of label two under addition modulo three.

\paragraph{Solution.}
An inverse label $a$ satisfies $2+a=0$ modulo three. The label one works
because $2+1=3$, congruent to zero. Geometrically a $120$-degree turn
undoes a $240$-degree turn up to one full revolution.

\paragraph{2. Check a bracket combination.}
Compute $[H,E+F]$ for the matrices above.

\paragraph{Solution.}
Bilinearity and the computed relations give
$[H,E+F]=[H,E]+[H,F]=2E-2F$.
Thus $E+F$ is not an eigenvector of this map: its two components acquire
opposite factors, and $2E-2F$ is not a scalar multiple of $E+F$.

\paragraph{3. Reflect an $A_2$ root.}
Find $s_{\alpha_1}(\alpha_2)$ and check its squared length.

\paragraph{Solution.}
Because $\alpha_2\cdot\alpha_1=-1$ and $\alpha_1^2=2$,
$s_{\alpha_1}(\alpha_2)=\alpha_2+\alpha_1=e_1-e_3$.
Its squared length is two, so the result is another root of the same length.
Chapter~\ref{learn:advanced} extends the distinction between algebraic
structure and physical interpretation.

\chapter{Geometry and analysis: distance, dimension, and complex structure}
\label{learn:geometry}

\begin{learningcard}{Prerequisites and scope}
Bring real limits and derivatives from Chapter~\ref{learn:calculus}, partial
derivatives from Chapter~\ref{learn:multivariable}, and complex arithmetic
from precalculus. Three questions guide this chapter: what does ``nearby''
mean, how can a set have fractional dimension, and what extra structure makes
a complex function modular? These are entrances to topology, fractal geometry,
and complex analysis, with a worked example at each entrance.
\end{learningcard}

\section{Distance comes with axioms}
A metric on a set $X$ is a function $d:X\times X\to\RR$ satisfying
$d(x,y)\geq0$, equality exactly when $x=y$, symmetry
$d(x,y)=d(y,x)$, and the triangle inequality
$d(x,z)\leq d(x,y)+d(y,z)$. Euclidean distance is an example. On $\RR^2$,
the Manhattan distance $|x_1-y_1|+|x_2-y_2|$ is another, useful when
horizontal and vertical travel supply the permitted paths.

The open ball $B_r(x)$ consists of points whose distance from $x$ is less
than $r$. A subset is open when every one of its points has an open ball
contained in the subset. Open sets record which perturbations stay inside
a region. A topology generalizes these neighborhood relations through a
collection of open sets, even when a preferred metric is absent.

A set is compact when every cover by open sets admits a finite subcover.
In Euclidean space, compactness is equivalent to being closed and bounded.
This equivalence is special to finite-dimensional Euclidean space and related
settings, rather than a definition valid in every metric space. A continuous
real function on a nonempty compact set attains its maximum and minimum.
This is the analysis input used in the symmetric spectral theorem of
Chapter~\ref{learn:linear_algebra}.

\section{Worked example: why an excluded endpoint matters}
On the closed interval $[0,1]$, the continuous function $f(x)=x$ attains
its maximum one. On the open interval $(0,1)$, it is bounded above by one
but has no maximum: for every proposed maximizing point $x$, the point
$(x+1)/2$ lies in the interval and has a larger function value.
The same formula has different extremum behavior because the domain changed.

The interval $(0,1)$ also has the open cover
$\{(1/n,1):n\geq2\}$, interpreted as open subsets of $(0,1)$.
Every point is covered by choosing $n>1/x$. A finite selection has a
largest index $N$ and leaves points at or below $1/N$ uncovered.
The failure of compactness is visible through the definition itself.

\section{Worked example: count scales in the Cantor set}
Start with $[0,1]$, remove the open middle third, then repeat the removal
inside every surviving interval. The middle-thirds Cantor set $C$ is the
intersection of all stages. Stage $n$ contains $2^n$ closed intervals of
length $3^{-n}$. If $N(\epsilon)$ is the least number of intervals of
length at most $\epsilon$ needed to cover $C$, then at these scales
$N(3^{-n})=2^n$. The stage intervals supply a cover; selecting their
left endpoints supplies $2^n$ points mutually farther apart than $3^{-n}$,
so a shorter cover needs at least that many intervals.

\begin{figure}[htbp]
  \centering
  \begin{tikzpicture}[x=9cm,y=0.95cm,font=\small,line width=0.8pt]
  \foreach \stage in {0,1,2,3} {
    \node[left] at (-0.03,-\stage) {$n=\stage$};
  }
  \draw[line width=3pt] (0,0)--(1,0);
  \foreach \numerator in {0,2} {
    \draw[line width=3pt] ({\numerator/3},-1)--({(\numerator+1)/3},-1);
  }
  \foreach \numerator in {0,2,6,8} {
    \draw[line width=3pt] ({\numerator/9},-2)--({(\numerator+1)/9},-2);
  }
  \foreach \numerator in {0,2,6,8,18,20,24,26} {
    \draw[line width=3pt] ({\numerator/27},-3)--({(\numerator+1)/27},-3);
  }
  \node[above] at (0,0.12) {$0$};
  \node[above] at (1,0.12) {$1$};
  \node[right] at (1.04,0) {$N=1$, $\ell=1$};
  \node[right] at (1.04,-1) {$N=2$, $\ell=1/3$};
  \node[right] at (1.04,-2) {$N=4$, $\ell=1/9$};
  \node[right] at (1.04,-3) {$N=8$, $\ell=1/27$};
  \node[below] at (0.5,-3.4) {At stage $n$: $N=2^n$ intervals, each of length $\ell=3^{-n}$.};
\end{tikzpicture}

  \caption{Read downward from the whole interval to successive removals. Each surviving interval produces two intervals one third as long. The right-hand labels record the count and individual length whose logarithmic ratio gives the dimension.}
  \label{fig:learn:cantor_construction}
\end{figure}

When the limit exists, box dimension is
\[
 \dim_B C=\lim_{\epsilon\to0}
 \frac{\log N(\epsilon)}{\log(1/\epsilon)}.
\]
At stage scales the ratio is $\log2/\log3$. For intermediate scales,
monotonicity traps covering counts between neighboring stage counts, giving
the same limit. Thus $\dim_B C=\log2/\log3\approx0.63093$.

Hausdorff dimension instead permits covers with different small diameters.
For $s\geq0$, minimize $\sum_j(\operatorname{diam}U_j)^s$ over countable
covers with diameters at most $\delta$, then let $\delta\to0$ to define
Hausdorff $s$-measure. The dimension is the threshold above which that
measure is zero. The Cantor stage cover has cost
$2^n(3^{-n})^s=(2\,3^{-s})^n$, proving the upper bound
$\dim_H C\leq\log2/\log3$.

For the matching lower bound, assign mass $2^{-n}$ to each stage-$n$
interval. These consistent assignments define the uniform Cantor probability
measure. An interval of length $\ell$, where
$3^{-n}\leq\ell<3^{-(n-1)}$, meets at most four stage-$n$ intervals.
Its mass is therefore at most $4\,2^{-n}\leq4\ell^s$ for
$s=\log2/\log3$. Any small cover of $C$ consequently has
$1\leq4\sum_j(\operatorname{diam}U_j)^s$, by containing each covering
set in an interval of its diameter. The Hausdorff $s$-measure is positive,
proving the lower bound. The two dimensions agree for this set; other sets
can separate them \citep{falconer2004fractal}.

\begin{learningcard}{Historical situation: a measurement question changes with scale}
Mandelbrot's 1967 paper asks how long Britain's coast is and relates
scale-dependent length measurements to statistical self-similarity and
fractional dimension \citep{mandelbrot1967long}. The publication supplies a
specific historical example of a measurement problem motivating geometric
language. An actual coastline has material and observational scale limits;
the exact Cantor construction is a mathematical model with indefinitely
repeated structure. A finite log--log fit to data estimates behavior over
its fitted scales rather than proving an asymptotic dimension.
\end{learningcard}

\section{Complex differentiation controls every approach direction}
For $z=x+iy$, a complex derivative is the limit
$f'(z)=\lim_{h\to0}[f(z+h)-f(z)]/h$ with complex $h$ approaching zero
from all directions. Write $f=u+iv$. Approaching along the real and imaginary
axes gives the necessary Cauchy--Riemann equations
\[
 u_x=v_y,\qquad u_y=-v_x.
\]
If the real partial derivatives are continuous in a neighborhood and satisfy
these equations there, the function is complex differentiable there. A
function complex differentiable throughout an open region is holomorphic.

For $f(z)=z^2$, $u=x^2-y^2$ and $v=2xy$ satisfy the equations, and
expanding $(z+h)^2$ gives $f'(z)=2z$. For $f(z)=\overline z$,
the real-axis difference quotient is one and the imaginary-axis quotient
is minus one, so the complex derivative fails everywhere. A real
coordinate formula alone therefore says little about holomorphy.

\section{Modular forms need a boundary condition too}
Let $\mathcal H=\{\tau\in\CC:\operatorname{Im}\tau>0\}$.
Integer matrices with determinant one form $\mathrm{SL}(2,\ZZ)$ and act by
$\tau\mapsto(a\tau+b)/(c\tau+d)$. Direct algebra gives
\[
 \operatorname{Im}\frac{a\tau+b}{c\tau+d}
 =\frac{\operatorname{Im}\tau}{|c\tau+d|^2}>0.
\]
Two basic transformations are $T(\tau)=\tau+1$ and $S(\tau)=-1/\tau$.
At $\tau=i$, $S(i)=i$ while $T(i)=1+i$; both stay in $\mathcal H$.

A modular form of integer weight $k$ for the full modular group is a
holomorphic function on $\mathcal H$ satisfying
$f((a\tau+b)/(c\tau+d))=(c\tau+d)^k f(\tau)$ and holomorphic at
the cusp. For this group the cusp condition means that its expansion in
$q=e^{2\pi i\tau}$ near $q=0$ has only nonnegative powers.
For subgroups, every inequivalent cusp must be checked. A cusp form has
zero constant term at each cusp. These conditions distinguish a modular
form from a function that merely obeys a transformation law
\citep{diamond2005first}. The invariant $j$, whose expansion starts with
$q^{-1}$, has a pole at the cusp and is a modular function, not a
holomorphic modular form there. Classification and coefficient theorems
require a further course in complex analysis and modular forms.

\section{Exercises with full solutions}
\paragraph{1. Compare two metrics.}
Find the Euclidean and Manhattan distances between $(0,0)$ and $(3,4)$.

\paragraph{Solution.}
Euclidean distance is $\sqrt{3^2+4^2}=5$. Manhattan distance is
$|3|+|4|=7$. Both obey metric axioms, but they encode different path costs.
The pair of values demonstrates that choosing a metric is substantive.

\paragraph{2. Recover a scaling exponent.}
A separated self-similar construction keeps four pieces, each scaled by
$1/5$. Find its similarity dimension, and state the needed qualification.

\paragraph{Solution.}
The scaling equation is $4(1/5)^s=1$, giving $s=\log4/\log5$.
For similarities satisfying an appropriate separation condition, such as
the open set condition, this exponent equals Hausdorff dimension. Without
such a condition, overlapping pieces can invalidate the inference.

\paragraph{3. Use a modular transformation at a fixed point.}
If $f$ has weight two for the full modular group, what does $S(i)=i$
force about $f(i)$?

\paragraph{Solution.}
The transformation law for $S$ gives $f(-1/\tau)=\tau^2 f(\tau)$.
At $i$, it becomes $f(i)=i^2f(i)=-f(i)$, hence $f(i)=0$.
This is a necessary consequence of the transformation law; it alone
constructs neither a nonzero form nor its coefficients.

\chapter{Entering the advanced album: algebra, quantum states, and fluids}
\label{learn:advanced}

\begin{learningcard}{Prerequisites and three destinations}
Use vector spaces and inner products, probability, algebraic structure, and
numerical conservation from the preceding chapters. We now construct one
bounded entry example for each of three advanced routes: Cayley--Dickson
algebras, quantum states, and lattice-Boltzmann computation. The examples
supply enough vocabulary and checks to read the advanced review critically;
they leave specialized classification, quantum information, and fluid
analysis as further coursework.
\end{learningcard}

\section{Doubling changes multiplication as well as dimension}
\label{learn:hypercomplex}
Starting with $\RR$ and its identity conjugation, the Cayley--Dickson
construction replaces an algebra by ordered pairs. Fix the convention
\[
 (a,b)(c,d)=(ac-\overline d\,b,\;da+b\overline c),
 \qquad \overline{(a,b)}=(\overline a,-b).
\]
The real vector-space dimension doubles at each step. Basis vectors are
ordered recursively: the old basis occupies the first half, followed by
pairs with zero first component and old basis in the second half. With
$e_0=1$, this fixes every multiplication sign used below.
Successive dimensions one, two, four, and eight give
$\RR,\CC,\HH,\OO$. The product loses commutativity at quaternions and
associativity at octonions. Norm multiplicativity survives through octonions
but fails for sedenions in dimension sixteen \citep{baez2001octonions}.

\section{Worked example: quaternion arithmetic and a boundary}
Write a quaternion as $q=a+bi+cj+dk$, with
$i^2=j^2=k^2=-1$, $ij=k$, $jk=i$, $ki=j$, and reversed distinct products
having the opposite sign. Then
\[
 (1+i)(1+j)=1+i+j+k,\qquad
 (1+j)(1+i)=1+i+j-k.
\]
Conjugation sends $q$ to $\overline q=a-bi-cj-dk$, and
$q\overline q=a^2+b^2+c^2+d^2=\|q\|^2$. For $q\ne0$,
$q^{-1}=\overline q/\|q\|^2$. Associativity and
$\overline{pq}=\overline q\,\overline p$ give
\[
 \|pq\|^2=pq\overline q\,\overline p
 =\|q\|^2p\overline p=\|p\|^2\|q\|^2.
\]
The scalar interpretation of the products justifies the norm equation.
For the displayed example, each factor has norm $\sqrt2$ and the product
has norm two. Noncommutativity coexists with a multiplicative norm.

A zero divisor is a nonzero element whose product with another nonzero
element is zero. Under the fixed doubling convention, the sedenion vectors
$a=e_3+e_{10}$ and $b=e_6-e_{15}$ satisfy
\[
 e_3e_6=e_5,\quad e_3e_{15}=e_{12},\quad
 e_{10}e_6=e_{12},\quad e_{10}e_{15}=e_5,
 \qquad ab=e_5-e_{12}+e_{12}-e_5=0.
\]
Each factor has Euclidean norm $\sqrt2$, demonstrating norm failure as well
as zero divisors. The coordinate convention is essential to replaying this
witness. It also explains why extending a successful quaternion construction
to every doubled algebra requires fresh proofs of the properties used.

\begin{learningcard}{Historical situation: reconstruct a mathematical route}
Baez's survey of the octonions describes the progression through normed
division algebras and their historical development \citep{baez2001octonions}.
It is a modern scholarly account rather than a contemporaneous record of
an inventor's thoughts. Our doubling sequence is an instructional
reconstruction: it makes the algebraic dependencies visible without
asserting that every historical participant followed this exact path.
\end{learningcard}

\section{Worked example: probabilities from a complex state}
\label{learn:quantum}
A pure qubit is represented by a unit vector
$\psi=\alpha|0\rangle+\beta|1\rangle$ in $\CC^2$, with
$|0\rangle=(1,0)^T$, $|1\rangle=(0,1)^T$ and
$|\alpha|^2+|\beta|^2=1$. Vectors differing by a common phase
$e^{i\theta}$ represent the same pure state. The complex inner product is
$\langle\phi,\psi\rangle=\sum_j\overline{\phi_j}\psi_j$.
The Born rule assigns computational-basis outcome probabilities
$P(0)=|\alpha|^2$ and $P(1)=|\beta|^2$.

A unitary matrix $U$ satisfies $U^\dagger U=I$, where the dagger denotes
conjugate transpose. It preserves inner products and normalization because
$\langle U\phi,U\psi\rangle=\phi^\dagger U^\dagger U\psi
=\langle\phi,\psi\rangle$. For the Hadamard matrix
\[
 H=\frac1{\sqrt2}\begin{pmatrix}1&1\\1&-1\end{pmatrix},\qquad
 H|0\rangle=\frac{|0\rangle+|1\rangle}{\sqrt2},
\]
the two measurement probabilities are $1/2$. Applying $H$ again gives
$H^2|0\rangle=|0\rangle$, so the final probability of zero is one.
Amplitudes combine before their magnitudes are squared. Replacing the
intermediate state by a classical random bit would erase relative-phase
information relevant to later operations.

The state-space postulates and Born rule are physical modeling rules, while
normalization preservation is a mathematical consequence of unitarity.
Connecting either to an experiment requires a specified preparation, operation,
and measurement. A classical program holding complex arrays does not by
itself demonstrate a quantum device. For the formal introductory framework,
see Nielsen and Chuang, Chapter 2.\footnote{Michael A. Nielsen and Isaac L.
Chuang, \emph{Quantum Computation and Quantum Information}, 10th anniversary
edition (2010), Chapter 2, especially the postulates of quantum mechanics:
\url{https://doi.org/10.1017/CBO9780511976667}.}

\section{A lattice-Boltzmann step has exact moment obligations}
\label{learn:fluids}
A lattice-Boltzmann method evolves populations $f_i$ associated with a finite
set of velocities $c_i$. In lattice units, define density and momentum by
$\rho=\sum_i f_i$ and $\rho u=\sum_i c_i f_i$ for $\rho>0$.
For the square D2Q9 stencil, velocities are $(0,0)$, the four axial unit
vectors, and the four diagonal vectors $(\pm1,\pm1)$. Their weights are
$4/9$, $1/9$ per axial vector, and $1/36$ per diagonal vector.
The weights sum to one; symmetry gives
\[
 \sum_iw_ic_i=0,\qquad
 \sum_iw_i c_{i\alpha}c_{i\beta}=\frac13\delta_{\alpha\beta},
\]
where $\delta_{\alpha\beta}$ is one for equal indices and zero otherwise.
Odd third moments vanish by opposite-velocity pairing.

The common low-velocity equilibrium is
\[
 f_i^{\mathrm{eq}}=w_i\rho\left[
 1+3c_i\cdot u+\frac92(c_i\cdot u)^2-\frac32|u|^2\right].
\]
Summing the polynomial with the weight identities gives
$\sum_i f_i^{\mathrm{eq}}=\rho[1+(9/2)(|u|^2/3)-(3/2)|u|^2]=\rho$.
Multiplying by $c_i$ and summing gives $\sum_i c_if_i^{\mathrm{eq}}=\rho u$:
the linear term contributes $3\rho u/3$, and the other vector sums vanish.
Thus the collision step
$f_i^*=f_i-\omega(f_i-f_i^{\mathrm{eq}})$ preserves both moments in exact
arithmetic. Streaming $f_i(x+c_i,t+1)=f_i^*(x,t)$ permutes populations on a
periodic lattice, so global mass is preserved as well.

Deriving incompressible Navier--Stokes behavior requires further assumptions:
low Mach number, suitable scale separation, the appropriate higher stencil
moments, and compatible boundary and forcing rules. In the standard BGK
hydrodynamic limit, lattice-unit kinematic viscosity is
$\nu=(1/3)(1/\omega-1/2)$; $0<\omega<2$ makes that coefficient positive,
but is insufficient alone to guarantee every simulation's stability.
Diffusion equations can use related scalar-population constructions, with
different equilibrium and moment targets. Naming both approaches
``lattice-Boltzmann'' leaves their modeled quantities distinct.
For a derivation and its assumptions, consult the method review by Chen and
Doolen.\footnote{Shiyi Chen and Gary D. Doolen, ``Lattice Boltzmann Method
for Fluid Flows,'' \emph{Annual Review of Fluid Mechanics} 30 (1998),
pp.~329--364: \url{https://doi.org/10.1146/annurev.fluid.30.1.329}.}

\section{Exercises with full solutions}
\paragraph{1. Invert a quaternion.}
Find the inverse of $q=1+2i$ and check the product.

\paragraph{Solution.}
The squared norm is five, so $q^{-1}=(1-2i)/5$.
Multiplication gives $(1+2i)(1-2i)/5=(1-4i^2)/5=1$.
The reversed product also equals one. This calculation uses the quaternion
relations and proves an inverse for the specified element.

\paragraph{2. Keep relative phase until the measurement.}
Apply $H$ to $(|0\rangle-|1\rangle)/\sqrt2$ and find the outcome probabilities.

\paragraph{Solution.}
Matrix multiplication gives
$\frac12(1-1,1+1)^T=(0,1)^T=|1\rangle$.
The probabilities are zero for outcome zero and one for outcome one.
Before the transformation, both probabilities were $1/2$. The relative
minus sign changes the subsequent interference.

\paragraph{3. Detect a population defect.}
Suppose a faulty equilibrium sums to $\rho+\varepsilon$ at one lattice
site. What does collision do to that site's density?

\paragraph{Solution.}
Summing the update gives
$\sum_i f_i^*=\rho-\omega[\rho-(\rho+\varepsilon)]
=\rho+\omega\varepsilon$.
The defect injects density in proportion to the relaxation factor. Periodic
streaming redistributes that injected amount and preserves its global sum;
it cannot repair the collision defect. This is why moment identities should
be checked before interpreting pictures of a flow.

\section{Read the existing advanced review by its evidence contract}
The advanced mathematical baseline studies exceptional Lie algebras,
Cayley--Dickson properties, modular forms, and fractal dimensions. The
computational portfolio then asks what retained artifacts establish about
implementations and proposed physical mechanisms. Carry three questions into
each entry: which mathematical statement has a proof, which finite calculation
has been reproduced, and which physical inference needs an independently
specified model and observation? Subsequent courses can deepen representation
theory, complex analysis, functional analysis, quantum information, and
continuum mechanics along separate paths. These chapters establish a shared
entry language and worked checks, rather than completion of those subjects.

\appendix
\chapter{Readiness solutions}

\label{learn:solutions}

Compare the operations as well as the final value. Return to the named lesson when a step needs repair.

\section{Solutions: From convincing examples to proof}

\label{solution:entry:proof}

\textbf{Entry answer.} Distribute multiplication to obtain x\textasciicircum{}2+2x+1.\par

\label{solution:ready:proof}

\textbf{Readiness answer.} Write the integers as 2m+1 and 2n+1 with integer m,n. Their sum is 2(m+n+1), twice an integer.\par

\hyperref[learn:proof]{Return to the lesson} or \hyperref[readiness:proof]{return to its next-question check}.

\section{Solutions: From balanced equations to matrix maps}

\label{solution:entry:matrix_maps}

\textbf{Entry answer.} Subtract the first equation from the second: x=4. Substitute to get y=3.\par

\label{solution:ready:matrix_maps}

\textbf{Readiness answer.} The outputs are 11 and 12. Their difference is -1, also obtained from the composite row (2,-1) applied to (2,5).\par

\hyperref[learn:matrix_maps]{Return to the lesson} or \hyperref[readiness:matrix_maps]{return to its next-question check}.

\section{Solutions: Complex numbers: multiplication, length, and conjugation}

\label{solution:entry:complex_numbers}

\textbf{Entry answer.} Distribution gives ac+ad+bc+bd; complex multiplication also uses i\textasciicircum{}2=-1.\par

\label{solution:ready:complex_numbers}

\textbf{Readiness answer.} The conjugate is 3-4i, squared modulus is 25, and reciprocal is (3-4i)/25. Multiplication checks a product of one.\par

\hyperref[learn:complex_numbers]{Return to the lesson} or \hyperref[readiness:complex_numbers]{return to its next-question check}.

\section{Solutions: Linear algebra: coordinates, maps, and modes}

\label{solution:entry:linear_algebra}

\textbf{Entry answer.} The output is (1,-1).\par

\label{solution:ready:linear_algebra}

\textbf{Readiness answer.} A(1,1)=3(1,1) and A(1,-1)=(1,-1). The two vectors are independent because neither is a scalar multiple of the other.\par

\hyperref[learn:linear_algebra]{Return to the lesson} or \hyperref[readiness:linear_algebra]{return to its next-question check}.

\section{Solutions: Calculus: local change and accumulated change}

\label{solution:entry:calculus}

\textbf{Entry answer.} The geometric tail is 2\textasciicircum{}(1-n), tending to zero as n increases.\par

\label{solution:ready:calculus}

\textbf{Readiness answer.} Subtract x\textasciicircum{}2 from (x+h)\textasciicircum{}2 and divide by h to obtain 2x+h. Let h approach zero to obtain 2x.\par

\hyperref[learn:calculus]{Return to the lesson} or \hyperref[readiness:calculus]{return to its next-question check}.

\section{Solutions: Several variables: local maps and conservation}

\label{solution:entry:multivariable}

\textbf{Entry answer.} The derivative is 2x; the dot product is 11.\par

\label{solution:ready:multivariable}

\textbf{Readiness answer.} The partial derivatives are 2x and 6y, giving gradient (2,12). Its dot product with a unit direction gives the directional rate.\par

\hyperref[learn:multivariable]{Return to the lesson} or \hyperref[readiness:multivariable]{return to its next-question check}.

\section{Solutions: Dynamics: equations, modes, and trustworthy computation}

\label{solution:entry:dynamics}

\textbf{Entry answer.} The chain rule gives -2 exp(-2t).\par

\label{solution:ready:dynamics}

\textbf{Readiness answer.} Solve -1\textless{}1-2h\textless{}1 to obtain 0\textless{}h\textless{}1. At h=1 the factor is -1 and the numerical solution does not decay.\par

\hyperref[learn:dynamics]{Return to the lesson} or \hyperref[readiness:dynamics]{return to its next-question check}.

\section{Solutions: Probability: uncertainty, evidence, and comparison}

\label{solution:entry:probability}

\textbf{Entry answer.} The fraction is 4/100=1/25.\par

\label{solution:ready:probability}

\textbf{Readiness answer.} There are 90+495=585 positives. The posterior is 90/585=2/13. The denominator counts observed positives, rather than the whole population.\par

\hyperref[learn:probability]{Return to the lesson} or \hyperref[readiness:probability]{return to its next-question check}.

\section{Solutions: Numerical reasoning: approximation with an error budget}

\label{solution:entry:numerical}

\textbf{Entry answer.} The discrete factor changes with h; its repeated product approximates the exponential only in a controlled limit.\par

\label{solution:ready:numerical}

\textbf{Readiness answer.} One panel gives 1/2 with error 1/6. Two give 3/8 with error 1/24, a factor-four reduction for this exact example.\par

\hyperref[learn:numerical]{Return to the lesson} or \hyperref[readiness:numerical]{return to its next-question check}.

\section{Solutions: Algebraic structure: symmetry, brackets, and roots}

\label{solution:entry:algebra}

\textbf{Entry answer.} The first basis vector selects the first column and the second selects the second.\par

\label{solution:ready:algebra}

\textbf{Readiness answer.} AB is diag(1,0), BA is diag(0,1), so the commutator is diag(1,-1). The nonzero result records order dependence.\par

\hyperref[learn:algebra]{Return to the lesson} or \hyperref[readiness:algebra]{return to its next-question check}.

\section{Solutions: Geometry and analysis: distance, dimension, and complex structure}

\label{solution:entry:geometry}

\textbf{Entry answer.} The modulus is 5, obtained from the square root of 3\textasciicircum{}2+4\textasciicircum{}2.\par

\label{solution:ready:geometry}

\textbf{Readiness answer.} There are 2\textasciicircum{}n intervals of length 3\textasciicircum{}(-n). The self-similar dimension solves 2*(1/3)\textasciicircum{}d=1, so d=log(2)/log(3). The chapter explains the additional dimension argument.\par

\hyperref[learn:geometry]{Return to the lesson} or \hyperref[readiness:geometry]{return to its next-question check}.

\section{Solutions: Doubling changes multiplication as well as dimension}

\label{solution:entry:hypercomplex}

\textbf{Entry answer.} The cross terms cancel and i\textasciicircum{}2=-1, giving 2.\par

\label{solution:ready:hypercomplex}

\textbf{Readiness answer.} In a division algebra, left multiplication by any nonzero element is invertible and hence injective. A nonzero element mapped to zero violates injectivity.\par

\hyperref[learn:hypercomplex]{Return to the lesson} or \hyperref[readiness:hypercomplex]{return to its next-question check}.

\section{Solutions: Worked example: probabilities from a complex state}

\label{solution:entry:quantum}

\textbf{Entry answer.} Conjugate-first multiplication gives 1+1=2.\par

\label{solution:ready:quantum}

\textbf{Readiness answer.} Divide by sqrt(2). Each squared modulus is 1/2, summing to one. Interpreting those numbers as measurement probabilities uses the Born postulate.\par

\hyperref[learn:quantum]{Return to the lesson} or \hyperref[readiness:quantum]{return to its next-question check}.

\section{Solutions: A lattice-Boltzmann step has exact moment obligations}

\label{solution:entry:fluids}

\textbf{Entry answer.} Roundoff comes from finite arithmetic; truncation comes from replacing a limiting operation by a finite approximation.\par

\label{solution:ready:fluids}

\textbf{Readiness answer.} The sum is 4/9+4/9+1/9=1. Opposite velocities have equal weights, so each component of the first moment cancels.\par

\hyperref[learn:fluids]{Return to the lesson} or \hyperref[readiness:fluids]{return to its next-question check}.

\backmatter
\bibliographystyle{plainnat}
\bibliography{papers/references}
\end{document}
